\section{Wasserstein Gradient Flows}
\label{sec-wasserstein-gradient-flows}
\subsection{Minimizing Movements and Wasserstein Gradients}
\begin{equation}
\alpha_{t+\tau} \in \uargmin{\alpha\in\Pp_2(\RR^d)}
\left\{\frac{1}{2 \tau} \Wass_2(\alpha_t, \alpha)^2 + f(\alpha)\right\}. \label{eq:jko-discr}
\end{equation}
\paragraph{Euclidean gradient flows.}
If we restrict \eqref{eq:jko-discr} to finite dimensions and assume $\alpha_t = \delta_{x(t)}$ and $\alpha = \delta_x$ (single Dirac measures), this matches the implicit Euler scheme:

\begin{equation*}
x(t+\tau) \in \uargmin{x}\left\{\frac{1}{2 \tau} \norm{x - x(t)}^2 + h(x)\right\},
\end{equation*}
where $h(x) = f(\delta_x)$. When $h$ is differentiable and the minimizer is unique, the solution is characterized by the implicit Euler formula

\(x(t+\tau) = (\Id + \tau \nabla h)^{-1}(x(t)).\)
\(x(t+\tau) = (\Id - \tau \nabla h)(x(t)) = x(t) - \tau \nabla h(x(t)).\)
\begin{equation}
\dot{x}(t) = -\nabla h(x(t)). \label{eq:grad-flow-classical}
\end{equation}
\paragraph{Wasserstein gradient formula.}
\(f((1-\tau)\alpha + \tau \beta) = f(\alpha + \tau \eta) = f(\alpha) + \tau \int [\delta f(\alpha)](x) \d \eta(x) + o(\tau).\)
\begin{defn}[Wasserstein gradient]\label{def-wasserstein-gradient}
Assume that $f$ admits a smooth first variation $\delta f(\alpha)$. In the smooth formal calculus on $\Pp_2(\RR^d)$, the Wasserstein gradient of $f$ at $\alpha$ is the gradient vector field
\(\Wgrad f(\alpha) = \nabla_x \delta f(\alpha).\)
\end{defn}
\begin{equation}
\frac{\partial \alpha_t}{\partial t} + \diverg(-\Wgrad f(\alpha_t) \alpha_t) = 0. \label{eq:wassflow-pde}
\end{equation}
\begin{proposition}[Formal Wasserstein gradient]\label{prop-formal-wass-gradient}
Assume that $f$ admits a smooth first variation $\delta f(\alpha)$ and that $\alpha$ has a smooth positive density. Let $v$ be a sufficiently regular velocity field, compactly supported or satisfying a no-flux boundary condition. For perturbations generated by $(\Id+\tau v)_\sharp\alpha$, the differential of $f$ is
\(\frac{\d}{\d\tau}_{|\tau=0} f\big((\Id+\tau v)_\sharp\alpha\big) = \int \dotp{\nabla \delta f(\alpha)(x)}{v(x)}\d\alpha(x).\)
Hence, for the Riemannian metric $\norm{v}_{L^2(\alpha)}^2=\int \norm{v}^2\d\alpha$, the Wasserstein gradient is the vector field
\(\Wgrad f(\alpha)=\nabla \delta f(\alpha).\)
\end{proposition}
\begin{proof}
The push-forward expansion gives, in the sense of distributions, \((\Id+\tau v)_\sharp\alpha = \alpha-\tau\diverg(\alpha v)+o(\tau).\)
Using the definition of the first variation, \(f((\Id+\tau v)_\sharp\alpha) = f(\alpha)-\tau\int \delta f(\alpha)\diverg(\alpha v)\d x+o(\tau).\)
An integration by parts, with either compact support or vanishing boundary flux, gives \(-\int \delta f(\alpha)\diverg(\alpha v)\d x = \int \dotp{\nabla\delta f(\alpha)}{v}\d\alpha.\)
By definition of the Riesz representative for the $L^2(\alpha)$ metric, this representative is $\nabla\delta f(\alpha)$.
\end{proof}
\paragraph{From the JKO step to the velocity field.}
A first-order expansion of the JKO step explains why~\eqref{eq:wassflow-pde} uses the vector field $\Wgrad f(\alpha)$. This is a formal local argument: assume that $\alpha_t$ is absolutely continuous and that, for the nearby competitors selected by the JKO step, the optimal map from $\alpha_t$ can be written as $\Id+\tau v$.

\(\Wass_2^2\bigl(\alpha_t,(\Id+\tau v)_\sharp\alpha_t\bigr) = \tau^2\norm{v}_{L^2(\alpha_t)}^2.\)
\(\min_v \frac{\tau}{2}\norm{v}_{L^2(\alpha_t)}^2 +f\bigl((\Id+\tau v)_\sharp\alpha_t\bigr),\)
\begin{align*}
(\Id+\tau v)_\sharp\alpha_t
&=
\alpha_t-\tau\diverg(v\alpha_t)+o(\tau),\\
f\bigl((\Id+\tau v)_\sharp\alpha_t\bigr)
&=
f(\alpha_t)
+\tau\int\dotp{\nabla_x\delta f(\alpha_t)(x)}{v(x)}\d\alpha_t(x)
+o(\tau)
\end{align*}
\[
\min_v f(\alpha_t)
+\tau\int\left[
\frac12\norm{v(x)}^2
+\dotp{\Wgrad f(\alpha_t)(x)}{v(x)}
\right]\d\alpha_t(x)
+o(\tau).
\]
\paragraph{Metric derivative and curves of maximal slope.}
The same descent principle admits a coordinate-free formulation, which is the right language for limits of JKO schemes. Let $(\mathcal X,d)$ be a metric space and let $x:[0,T]\to\mathcal X$ be an absolutely continuous curve.

\(|\dot x_t| \eqdef \lim_{h\to 0}\frac{d(x_{t+h},x_t)}{|h|}\)
\[
|\partial f|(x)
\eqdef
\limsup_{\substack{y\to x\\ y\neq x}}
\frac{( f(x)-f(y) )_+}{d(x,y)}.
\]
When this slope is a strong upper gradient for $f$, meaning that $|f(y_t)-f(y_s)|\leq \int_s^t |\partial f|(y_r)|\dot y_r|\d r$ along absolutely continuous curves $y_r$, an absolutely continuous curve $x_t$ is called a curve of maximal slope if it satisfies the energy-dissipation inequality

\(f(x_t) +\frac12\int_s^t |\dot x_r|^2\d r +\frac12\int_s^t |\partial f|^2(x_r)\d r \leq f(x_s), \qquad 0\leq s<t.\)
\paragraph{Discrete evolutions.}
If $f$ admits a first variation whose spatial gradient is well defined at empirical measures, the characteristic form of~\eqref{eq:wassflow-pde} preserves the number and weights of the atoms. For equal weights, write $\alpha_t = \frac{1}{n} \sum_i \delta_{x_i(t)}$. The particles $X(t):= (x_i(t))_i$ evolve according to the coupled ODEs

\begin{equation}
\dot{x}_i(t) = -n\nabla_{x_i} F(X(t)), \label{eq:wassflows-particles}
\end{equation}
where $F(X):= f\left(\frac{1}{n} \sum_i \delta_{x_i}\right)$ and the factor $n$ comes from the empirical Wasserstein metric $\frac1n\sum_i \norm{\dot x_i}^2$.

\paragraph{Linear functionals.}
\paragraph{Shannon neg-entropy.}
\paragraph{Interaction energies.}
\begin{equation}
f(\alpha):= \iint k(x, y) \d \alpha(x) \d \alpha(y). \label{eq:quadratic-func}
\end{equation}
For a symmetric kernel $k$:

\(\delta f(\alpha)(x) = 2 \int k(x, y) \d \alpha(y), \quad \Wgrad f(\alpha)(x) = 2 \int \nabla_x k(x, y) \d \alpha(y).\)
For $\alpha_0 = \frac{1}{n} \sum_i \delta_{x_i}$, the flow \eqref{eq:wassflow-pde} implies particles $(x_i(t))_i$ obey:

\(\dot{x}_i(t) = -\frac{2}{n} \sum_j \nabla k(x_i(t), x_j(t)).\)
If $k$ is positive definite, or more generally conditionally positive definite on signed measures of zero total mass as for the energy-distance kernel $k(x,y)=-\norm{x-y}$, and one minimizes the squared kernel discrepancy to a teacher distribution $\beta$, then

\[
\norm{\alpha-\beta}_k^2
=
\iint k\d\alpha\d\alpha
-2\int\!\left(\int k(x,y)\d\beta(y)\right)\d\alpha(x)
+\mathrm{constant}.
\]
Thus MMD-type training energies are exactly an interaction energy plus a linear potential; the teacher distribution appears through the potential $x\mapsto-2\int k(x,y)\d\beta(y)$. The corresponding empirical Wasserstein gradient flow is

\(\dot x_i(t) = -\frac{2}{n}\sum_j\nabla_x k(x_i(t),x_j(t)) + 2\int\nabla_x k(x_i(t),y)\d\beta(y).\)
\paragraph{Stochastic particles and McKean--Vlasov limits.}
\(\d X_t=b(X_t)\d t+\sqrt{2}\sigma\d B_t,\)
and the one-particle law $\alpha_t=\rho_t\d x$ directly satisfies the linear Fokker--Planck equation

\(\partial_t\rho_t=-\diverg(b\rho_t)+\sigma^2\Delta\rho_t.\)
\(\d X_i^n(t)=b(X_i^n(t),\al_t^n)\d t+\sqrt{2}\sigma\d B_i(t), \qquad \al_t^n=\frac1n\sum_{i=1}^n\delta_{X_i^n(t)}.\)
For finite $n$, the empirical law $\al_t^n$ is itself random. Under suitable Lipschitz, growth and chaotic-initialization assumptions, propagation of chaos states that finitely many particles become asymptotically independent as $n\to\infty$, all with the same deterministic law $\rho_t\d x$; equivalently, the empirical measure $\al_t^n$ converges in probability to this law.

\(\partial_t \rho_t = -\diverg\big(b(x,\rho_t)\rho_t\big) + \sigma^2\Delta\rho_t.\)
\(b(x,\rho)=-\nabla \frac{\delta \mathcal E}{\delta \rho}(x),\)
\(\mathcal E(\rho)+\sigma^2\int \rho\log\rho\,\d x.\)
For a prescribed target $\beta=\rho_\beta\d x$, take $b=\sigma^2\nabla\log\rho_\beta$. The SDE, the Fokker--Planck equation, and the deterministic continuity equation with score velocity are then three representations of the same gradient flow:

\begin{equation}\label{eq-relative-entropy-three-representations}
\partial_t\rho_t
=
\sigma^2\diverg\!\left(\rho_t\nabla\log\frac{\rho_t}{\rho_\beta}\right),
\qquad
v_t
=
\sigma^2\bigl(\nabla\log\rho_\beta-\nabla\log\rho_t\bigr).
\end{equation}
\paragraph{Gradient flows under density constraints.}
\begin{equation}\label{eq-density-constrained-energy}
f_{\rho_{\max}}(\rho\d x)=\int_\Omega h(x)\rho(x)\d x+
\iota_{\mathcal K_{\rho_{\max}}}(\rho\d x),
\qquad
\mathcal K_{\rho_{\max}}=
\enscond{\rho\d x\in\Pp_2(\Omega)}{0\leq\rho\leq\rho_{\max}\ \text{a.e.}},
\end{equation}
where $\Omega\subset\RR^d$ is a convex domain and $\mathcal K_{\rho_{\max}}$ is assumed nonempty; if $\Omega$ has finite volume, this requires $\rho_{\max}|\Omega|\geq1$. The indicator term has no classical first variation; it contributes the normal cone to $\mathcal K_{\rho_{\max}}$ in the Wasserstein first-order condition.

\begin{equation}\label{eq-density-constrained-jko}
\alpha^{k+1}
\in
\uargmin{\alpha\in\mathcal K_{\rho_{\max}}}
\frac{1}{2\tau}\Wass_2^2(\alpha^k,\alpha)+\int_\Omega h\,\d\alpha.
\end{equation}
\begin{equation}\label{eq-density-constrained-pressure-flow}
\partial_t\rho_t=\diverg\left(\rho_t\nabla(h+p_t)\right),
\qquad
0\leq\rho_t\leq\rho_{\max},
\qquad
p_t\geq0,
\qquad
p_t(\rho_{\max}-\rho_t)=0.
\end{equation}
The pressure vanishes in the unsaturated region and acts only where $\rho_t=\rho_{\max}$, pushing mass away from congested zones so that the density cap remains satisfied.

\paragraph{Multi-species gradient flows.}
Several applications involve several densities living on the same physical space: chemical species, color channels, cell populations or competing phases. Let $\Omega\subset\RR^d$ be the physical domain and assume that the component masses $\mathbf m=(m_1,\ldots,m_p)$ are positive and fixed.

\begin{equation}\label{eq-multispecies-space}
\Pp_{2,\mathbf m}(\Omega;\RR_+^p)=
\enscond{\alpha=(\alpha_1,\ldots,\alpha_p)}
{\alpha_i\in\Mm_+(\Omega),\ \alpha_i(\Omega)=m_i,\ \int_\Omega\norm{x}^2\d\alpha_i(x)<+\infty}.
\end{equation}
\begin{equation}\label{eq-multispecies-product-metric}
\Wass_{2,\oplus}^2(\alpha,\eta)=
\sum_{i=1}^p m_i\,
\Wass_2^2\!\left(\frac{\alpha_i}{m_i},\frac{\eta_i}{m_i}\right).
\end{equation}
\begin{equation}\label{eq-multispecies-jko}
\alpha^{k+1}
\in
\uargmin{\alpha\in\Pp_{2,\mathbf m}(\Omega;\RR_+^p)}
\frac{1}{2\tau}\Wass_{2,\oplus}^2(\alpha^k,\alpha)+f(\alpha).
\end{equation}
For smooth densities and first variations $\phi_i=\delta f/\delta\rho_i$, this yields

\begin{equation}\label{eq-multispecies-pde}
\partial_t\rho_i=\diverg\left(\rho_i\nabla\phi_i\right),
\qquad
\phi_i=\frac{\delta f}{\delta\rho_i},
\qquad i=1,\ldots,p.
\end{equation}
\begin{equation}\label{eq-multispecies-composition-constraint}
\mathcal C_\beta=
\enscond{\alpha\in\Pp_{2,\mathbf m}(\Omega;\RR_+^p)}
{\sum_{i=1}^p \alpha_i=\beta}.
\end{equation}
\begin{equation}\label{eq-multispecies-constrained-jko}
\alpha^{k+1}
\in
\uargmin{\alpha\in\mathcal C_\beta}
\frac{1}{2\tau}\Wass_{2,\oplus}^2(\alpha^k,\alpha)+f(\alpha).
\end{equation}
\begin{equation}\label{eq-multispecies-constrained-pde}
\partial_t\rho_i=\diverg\left(\rho_i\nabla(\phi_i-\lambda_t)\right),
\qquad
\diverg(b\nabla\lambda_t)=\diverg\left(\sum_{i=1}^p \rho_i\nabla\phi_i\right).
\end{equation}
For the separable Shannon entropy $f(\alpha)=\sum_i\int_\Omega \rho_i\log\rho_i\d x$, this becomes

\begin{equation}\label{eq-multispecies-modified-heat}
\partial_t\rho_i=\Delta\rho_i-\diverg(\rho_i\nabla\lambda_t),
\qquad
\diverg(b\nabla\lambda_t)=\Delta b.
\end{equation}
When $b$ is constant, the pressure can be chosen constant and the system reduces to independent heat equations which nevertheless preserve the pointwise sum $\sum_i\rho_i=b$. This product geometry should be contrasted with the vector-valued Benamou--Brenier distances of Section~\ref{sec-vector-matrix-valued-measures}, where the metric itself can couple the components through a non-diagonal mobility.

\subsection{Geodesic Convexity and Convergence}
\label{sec-geodesic-convexity}
\paragraph{Geodesics and convexity.}
\(\alpha_t=((1-t)P_0+tP_1)_\sharp\pi^\star, \qquad t\in[0,1],\)
where $P_0(x,y)=x$ and $P_1(x,y)=y$. If the optimal plan is induced by a Brenier map $T$, this reduces to $((1-t)\Id+tT)_\sharp\alpha_0$.

\begin{defn}[Geodesic convexity]\label{def-geodesic-convexity}
A functional $f$ on $\Pp_2(\RR^d)$ is geodesically convex if for every $\Wass_2$ geodesic $(\alpha_t)_t$, \(f(\alpha_t)\leq (1-t)f(\alpha_0)+t f(\alpha_1).\)
It is $\lambda$-geodesically convex if the right-hand side is improved by $-\frac{\lambda}{2}t(1-t)\Wass_2^2(\alpha_0,\alpha_1)$.
\end{defn}
\begin{prop}[Geodesic convexity of linear and quadratic energies]\phantomsection\label{prop-basic-geodesic-convexity}
\begin{enumerate}
\item If $h$ is convex, then the linear energy $\alpha\mapsto\int h\d\alpha$ is geodesically convex; if $h$ is $\lambda$-strongly convex, it is $\lambda$-geodesically convex. Conversely, if this linear energy is geodesically convex for all Dirac endpoints, then $h$ is convex.
\item More generally, let \(k\geq2\) and let \(H_k:(\RR^d)^k\to\RR\cup\{+\infty\}\) be convex. Then the polynomial energy
\(\alpha\mapsto \int_{(\RR^d)^k} H_k(x_1,\ldots,x_k)\,\d\alpha(x_1)\cdots\d\alpha(x_k)\)
is geodesically convex whenever the integral is well defined. If $H_k$ is $\lambda$-strongly convex on the product space, this energy is $k\lambda$-geodesically convex. In particular, the quadratic interaction
\(\alpha\mapsto\frac12\iint W(x,y)\d\alpha(x)\d\alpha(y)\)
is geodesically convex when \(W\) is convex on \(\RR^d\times\RR^d\), and it is $\lambda$-geodesically convex when $W$ is $\lambda$-strongly convex there.
\end{enumerate}
\end{prop}
\begin{proof}
Let \(\pi\) be an optimal coupling between \(\alpha_0\) and \(\alpha_1\), and set \(X_t=(1-t)X_0+tX_1\) under \((X_0,X_1)\sim\pi\). Convexity of \(h\) gives \(h(X_t)\leq(1-t)h(X_0)+t h(X_1)\), and strong convexity gives the additional term \(-\frac{\lambda}{2}t(1-t)\norm{X_0-X_1}^2\). Integrating proves the first claim. Necessity follows by testing on Dirac masses: the geodesic between \(\delta_x\) and \(\delta_y\) is \(\delta_{(1-t)x+ty}\), so geodesic convexity of \(\int h\d\alpha\) gives the usual convexity inequality for \(h\).

For the polynomial statement, take \(k\) independent copies \((X_0^\ell,X_1^\ell)_{\ell=1}^k\) of the same optimal coupling and define \(X_t^\ell=(1-t)X_0^\ell+tX_1^\ell\). Convexity of \(H_k\) on the product space gives
\(H_k(X_t^1,\ldots,X_t^k) \leq (1-t)H_k(X_0^1,\ldots,X_0^k)+tH_k(X_1^1,\ldots,X_1^k).\)
Integrating over the product coupling proves the claim. If $H_k$ is $\lambda$-strongly convex, the additional term is
\(-\frac{\lambda}{2}t(1-t) \sum_{\ell=1}^k\EE\norm{X_0^\ell-X_1^\ell}^2 = -\frac{k\lambda}{2}t(1-t)\Wass_2^2(\alpha_0,\alpha_1),\)
which proves the stated constant. For the quadratic interaction, the prefactor $1/2$ cancels the factor $k=2$, leaving the constant $\lambda$.
\end{proof}
\begin{prop}[One-dimensional interactions and energy-distance MMD]\label{prop-1d-interaction-halfline-convex}
Let \(\varphi:\RR_+\to\RR\) be continuous with at most quadratic growth, and define on \(\Pp_2(\RR)\)
\(\mathcal I_\varphi(\alpha) \eqdef \frac12\iint \varphi(|x-y|)\d\alpha(x)\d\alpha(y).\)
Then \(\mathcal I_\varphi\) is geodesically convex for \(\Wass_2\) on \(\Pp_2(\RR)\) if and only if \(\varphi\) is convex on \(\RR_+\).
Fix also \(\beta\in\Pp_2(\RR)\). For the conditionally positive kernel \(k(x,y)=-|x-y|\) from Definition~\ref{def-positive-kernels}, the squared MMD of Definition~\ref{def-kernel-mmd-norm}, equivalently the squared energy distance,
\(\mathcal E_\beta(\alpha) \eqdef \MMD_k^2(\alpha,\beta) = \iint -|x-y|\,\d(\alpha-\beta)(x)\d(\alpha-\beta)(y),\)
is geodesically convex as a function of \(\alpha\).
\end{prop}
\begin{proof}
Let \(Q_0,Q_1\) be the quantile functions of \(\alpha_0,\alpha_1\), and let \(Q_t=(1-t)Q_0+tQ_1\) be the quantile function of the one-dimensional \(\Wass_2\) geodesic \(\alpha_t\). For \(r>s\), monotonicity gives \(Q_i(r)-Q_i(s)\geq0\) for \(i=0,1\), and hence
\(Q_t(r)-Q_t(s) = (1-t)(Q_0(r)-Q_0(s))+t(Q_1(r)-Q_1(s)).\)
Using the symmetry of the double integral and the quantile parametrization, one can write
\(\mathcal I_\varphi(\alpha_t) = \int_0^1\int_0^r \varphi(Q_t(r)-Q_t(s))\,\d s\,\d r.\)
Convexity of \(\varphi\) on \(\RR_+\) gives the desired convexity after integration.

Conversely, test geodesic convexity on two-point measures \(\alpha_i=\frac12(\delta_0+\delta_{a_i})\), with \(a_i\geq0\). Their monotone geodesic is \(\alpha_t=\frac12(\delta_0+\delta_{(1-t)a_0+t a_1})\). Since
\(\mathcal I_\varphi(\alpha_t) = \frac14\varphi(0)+\frac14\varphi((1-t)a_0+t a_1),\)
and the identical \(\frac14\varphi(0)\) terms cancel from the geodesic-convexity inequality, one obtains the ordinary convexity inequality for \(\varphi\) on \(\RR_+\).

For the energy distance, the function \(\varphi(r)=-r\) is affine on \(\RR_+\), so the self-interaction \(-\iint |x-y|\d\alpha(x)\d\alpha(y)\) is affine along one-dimensional Wasserstein geodesics even though \(z\mapsto-|z|\) is not convex on \(\RR\). Write \(Q_\beta\) for the quantile function of \(\beta\). Expanding the MMD and dropping the term depending only on \(\beta\) gives
\[
\mathcal E_\beta(\alpha)
=
2\int_0^1\!\!\int_0^1 |Q_\alpha(r)-Q_\beta(s)|\,\d r\,\d s
-
\int_0^1\!\!\int_0^1 |Q_\alpha(r)-Q_\alpha(s)|\,\d r\,\d s
+\mathrm{constant}.
\]
Here the constant is independent of \(\alpha\). The positive cross-distance term is convex in \(Q_\alpha\), since the absolute value is convex and \(Q_t\) is affine in \(t\). The self-distance term is affine along quantile geodesics because
\(\int_0^1\!\!\int_0^1 |Q_t(r)-Q_t(s)|\,\d r\,\d s = 2\int_0^1\int_0^r (Q_t(r)-Q_t(s))\,\d s\,\d r.\)
Thus \(\alpha\mapsto\mathcal E_\beta(\alpha)\) is geodesically convex.
\end{proof}
\(\MMD_k^2(\alpha,\beta) = \iint k\,\d\alpha\d\alpha -2\iint k\,\d\alpha\d\beta +\mathrm{constant}\)
\begin{thm}[McCann displacement convexity for internal energies]\label{thm-mccann-internal-energy}
Let \(g:[0,+\infty)\to\RR\cup\{+\infty\}\) be convex with \(g(0)=0\), and define
\(\psi(r)=r^d g(r^{-d}),\qquad r>0.\)
If \(\psi\) is convex and nonincreasing on \((0,+\infty)\), then the internal energy
\[
\mathcal U_g(\alpha)=
\begin{cases}
\displaystyle \int_{\RR^d} g(\rho(x))\,\d x, & \alpha=\rho\d x,\\
+\infty, & \text{otherwise},
\end{cases}
\]
is geodesically convex on \(\Pp_2(\RR^d)\).
\end{thm}
\begin{proof}
It is enough to prove the inequality for smooth positive compactly supported densities; the general case follows by approximation and lower semicontinuity. Let \(T=\nabla\varphi\) be the Brenier map from \(\alpha_0=\rho_0\d x\) to \(\alpha_1=\rho_1\d x\), set \(T_t=(1-t)\Id+tT\), and denote
\(J_t(x)=\det\big((1-t)\Id+t\nabla T(x)\big).\)
The interpolation satisfies \(\alpha_t=(T_t)_\sharp\alpha_0\) and
\(\rho_t(T_t(x))J_t(x)=\rho_0(x).\)
By concavity of \(\det^{1/d}\) on positive semidefinite matrices, \(J_t(x)^{1/d}\geq (1-t)+t\det(\nabla T(x))^{1/d}.\)
Introduce \(s_t(x)=(J_t(x)/\rho_0(x))^{1/d}\). Since
\(\rho_1(T(x))\det\nabla T(x)=\rho_0(x)\), this gives
\(s_t(x)\geq (1-t)\rho_0(x)^{-1/d}+t\rho_1(T(x))^{-1/d}.\)
Using the change of variables \(y=T_t(x)\), one obtains
\(\int g(\rho_t(y))\,\d y =\int \rho_0(x)\,\psi(s_t(x))\,\d x.\)
Because \(\psi\) is nonincreasing and convex, \(\psi(s_t(x)) \leq (1-t)\psi(\rho_0(x)^{-1/d})+t\psi(\rho_1(T(x))^{-1/d}).\)
Multiplying by \(\rho_0\) and integrating gives
\(\mathcal U_g(\alpha_t)\leq(1-t)\mathcal U_g(\alpha_0)+t\mathcal U_g(\alpha_1)\).
\end{proof}
\begin{prop}[Geodesic convexity of power divergences]\label{prop-power-divergences-geodesic-convexity}\label{prop-kl-gibbs-geodesic-convexity}
For \(m>0\), define
\[
\phi_m(r)=
\begin{cases}
\dfrac{r^m-mr+m-1}{m(m-1)}, & m\neq1,\\[2mm]
r\log r-r+1, & m=1.
\end{cases}
\]
Then the following geodesic-convexity statements hold.
\begin{enumerate}
\item If \(\Omega\subset\RR^d\) is convex, \(\beta=b\,\mathds{1}_\Omega\d x\), and \(m\geq1-1/d\), then \(\alpha\mapsto\Divergm_{\phi_m}(\alpha|\beta)\) is geodesically convex on \(\Pp_2(\Omega)\).
\item If \(\beta=Z^{-1}e^{-V}\d x\) on \(\RR^d\), \(V\) is convex, and \(m\geq1\), then \(\alpha\mapsto\Divergm_{\phi_m}(\alpha|\beta)\) is geodesically convex on \((\Pp_2(\RR^d),\Wass_2)\).
\item In the KL case \(m=1\), if \(V\) is \(\lambda\)-strongly convex, then \(\alpha\mapsto\KL(\alpha|\beta)=\Divergm_{\phi_1}(\alpha|\beta)\) is \(\lambda\)-geodesically convex.
\end{enumerate}
All statements use the convention that the divergence is \(+\infty\) when \(\alpha\) is not absolutely continuous with respect to \(\beta\).
\end{prop}
\begin{proof}
For the flat reference case, write \(\alpha=\rho\d x\). Up to affine terms in \(\rho\), which only add constants on probability measures, the integrand \(b\,\phi_m(\rho/b)\) has non-affine part \(b^{1-m}\rho^m/(m(m-1))\) for \(m\neq1\), and \(\rho\log\rho\) for \(m=1\). For \(m=1\), McCann's criterion gives \(\psi(r)=-d\log r\). For \(m\neq1\), the transformed non-affine part is, up to the irrelevant factor \(b^{1-m}\),
\(\psi_m(r)=\frac{r^{d(1-m)}}{m(m-1)}.\)
If \(m>1\), this function is convex and nonincreasing. If \(0<m<1\), it is convex and nonincreasing exactly when \(d(1-m)\leq1\), i.e. \(m\geq1-1/d\). This proves the first claim by Theorem~\ref{thm-mccann-internal-energy}.

Assume next that \(m>1\) and \(V\) is convex. Up to constants, \(\Divergm_{\phi_m}(\alpha|\beta) = \frac{Z^{m-1}}{m(m-1)} \int \rho(x)^m e^{(m-1)V(x)}\,\d x.\)
Let \(T=\nabla\varphi\) be the Brenier map from \(\alpha_0=\rho_0\d x\) to \(\alpha_1\), set \(T_t=(1-t)\Id+tT\), and write
\(J_t(x)=\det((1-t)\Id+t\nabla T(x))\). Along the geodesic \(\alpha_t=(T_t)_\sharp\alpha_0\),
\[
\int \rho_t^m e^{(m-1)V}\d x
=
\int \rho_0(x)^m
\exp\!\left((m-1)\big(V(T_t(x))-\log J_t(x)\big)\right)\d x.
\]
For each \(x\), the function \(t\mapsto V(T_t(x))\) is convex, and \(t\mapsto-\log J_t(x)\) is convex because \(\log\det\) is concave on positive definite matrices. Since \(m-1>0\), the exponential of their sum is convex in \(t\). Integrating gives geodesic convexity. The nonsmooth case follows by approximation.

Finally, for \(m=1\), \(\KL(\alpha|\beta) = \int \rho\log\rho\,\d x+\int V\d\alpha+\log Z.\)
The entropy term is \(0\)-geodesically convex by Theorem~\ref{thm-mccann-internal-energy}. The linear potential term is \(\lambda\)-geodesically convex when \(V\) is \(\lambda\)-strongly convex, by Proposition~\ref{prop-basic-geodesic-convexity}. Taking \(\lambda=0\) also covers the \(m=1\) case in the second claim.
\end{proof}
\paragraph{Geodesically convex constraints.}
\label{par-geodesically-convex-constraints}
\begin{proposition}[Geodesic convexity of density caps]\label{prop-density-cap-geodesic-convex}
Assume that $\Omega\subset\RR^d$ is convex and let $\rho_{\max}>0$. The set $\mathcal K_{\rho_{\max}}$ in~\eqref{eq-density-constrained-energy} is geodesically convex in $\Pp_2(\Omega)$: if $\alpha_0,\alpha_1\in\mathcal K_{\rho_{\max}}$, then the constant-speed $\Wass_2$ geodesic $(\alpha_t)_{t\in[0,1]}$ between them also belongs to $\mathcal K_{\rho_{\max}}$.
\end{proposition}
\begin{proof}
Assume first that the endpoints have smooth positive densities $\alpha_i=\rho_i\d x$ with $\rho_i\leq\rho_{\max}$, and let $T=\nabla\varphi$ be the Brenier map from $\alpha_0$ to $\alpha_1$. Since $\Omega$ is convex, $T_t=(1-t)\Id+tT$ maps $\Omega$ into $\Omega$. The interpolation is $\alpha_t=(T_t)_\sharp\alpha_0$ and
\(\rho_t(T_t(x))=\frac{\rho_0(x)}{\det((1-t)\Id+t\nabla T(x))}.\)
The concavity of $\det^{1/d}$ on positive semidefinite matrices and the identity $\rho_1(T(x))\det\nabla T(x)=\rho_0(x)$ give
\(\rho_t(T_t(x))^{-1/d} \geq (1-t)\rho_0(x)^{-1/d}+t\rho_1(T(x))^{-1/d} \geq\rho_{\max}^{-1/d}.\)
Hence $\rho_t\leq\rho_{\max}$ for all $t$. The general case follows by approximation and lower semicontinuity of the $L^\infty$ bound, as in McCann's displacement-convexity theory.
\end{proof}
\paragraph{Dissipation and finite metric length.}
\(\frac{\d}{\d t}h(x_t) = -\norm{\nabla h(x_t)}^2 = -\norm{\dot x_t}^2.\)
Thus, for $0\leq s<t$, Cauchy's inequality gives

\begin{equation}\label{eq-euclidean-gradient-flow-length}
\int_s^t \norm{\dot x_r}\,\d r
\leq
\sqrt{t-s}\left(\int_s^t\norm{\dot x_r}^2\d r\right)^{1/2}
=
\sqrt{t-s}\,\bigl(h(x_s)-h(x_t)\bigr)^{1/2}.
\end{equation}
\begin{prop}[Finite metric length from dissipation]\label{prop-gradient-flow-finite-length}
Let $(\alpha_r)_{r\in[0,T]}$ be an absolutely continuous curve in $(\Pp_2(\RR^d),\Wass_2)$. Assume that, for all $0\leq s<t\leq T$, it satisfies the exact energy-dissipation identity
\(f(\alpha_s)-f(\alpha_t) = \int_s^t |\dot\alpha_r|_{\Wass_2}^2\d r.\)
Then
\begin{equation}\label{eq-wasserstein-gradient-flow-length}
\operatorname{Length}_{\Wass_2}\bigl((\alpha_r)_{r\in[s,t]}\bigr)
=
\int_s^t |\dot\alpha_r|_{\Wass_2}\,\d r
\leq
\sqrt{t-s}\,\bigl(f(\alpha_s)-f(\alpha_t)\bigr)^{1/2}.
\end{equation}
In the smooth Wasserstein gradient-flow setting with velocity $v_r=-\Wgrad f(\alpha_r)$, the dissipation identity reads equivalently
\(f(\alpha_s)-f(\alpha_t) = \int_s^t |\dot\alpha_r|_{\Wass_2}^2\d r = \int_s^t\int \norm{v_r(x)}^2\d\alpha_r(x)\d r.\)
If one assumes only the metric energy-dissipation inequality for a curve of maximal slope, then the same length estimate holds with the right-hand side multiplied by $\sqrt2$.
\end{prop}
\begin{proof}
Cauchy's inequality gives
\[
\int_s^t |\dot\alpha_r|_{\Wass_2}\,\d r
\leq
\sqrt{t-s}
\left(\int_s^t |\dot\alpha_r|_{\Wass_2}^2\d r\right)^{1/2}.
\]
The exact energy-dissipation identity substitutes the last integral by $f(\alpha_s)-f(\alpha_t)$, which proves~\eqref{eq-wasserstein-gradient-flow-length}. In the metric theory of curves of maximal slope, the energy-dissipation inequality gives instead
\(\frac12\int_s^t |\dot\alpha_r|_{\Wass_2}^2\d r \leq f(\alpha_s)-f(\alpha_t),\)
and the same argument yields the additional factor $\sqrt2$.
\end{proof}
\(\Wass_2(\alpha_s,\alpha_t) \leq \operatorname{Length}_{\Wass_2}\bigl((\alpha_r)_{r\in[s,t]}\bigr) \leq \sqrt{f(\alpha_0)-m_T}\,\sqrt{t-s},\)
with an additional factor \(\sqrt2\) under the energy-dissipation inequality. Thus an energy-dissipating trajectory is \(1/2\)-H\"older continuous as a curve in Wasserstein space on every time interval where the energy remains bounded from below.

\begin{prop}[Energy decay for convex Wasserstein flows]\label{prop-convex-wass-flow-rate}
Assume formally that \(f\) is geodesically convex, admits a smooth first variation, and has a minimizer \(\alpha^\star\). Let \((\alpha_t)_t\) be a smooth solution of the Wasserstein gradient flow
\(\partial_t\alpha_t+\diverg(\alpha_t v_t)=0, \qquad v_t=-\Wgrad f(\alpha_t).\)
Then
\(\frac{\d}{\d t} f(\alpha_t) = -\int \norm{\Wgrad f(\alpha_t)(x)}^2\,\d\alpha_t(x) \leq 0.\)
If \(T_t\) is the optimal map from \(\alpha_t\) to \(\alpha^\star\), then
\(f(\alpha_t)-f(\alpha^\star) \leq -\frac{\d}{\d t}\frac12 \Wass_2^2(\alpha_t,\alpha^\star),\)
and consequently
\(f(\alpha_t)-f(\alpha^\star) \leq \frac{\Wass_2^2(\alpha_0,\alpha^\star)}{2t}.\)
\end{prop}
\begin{proof}
The chain rule and Proposition~\ref{prop-formal-wass-gradient} give
\(\frac{\d}{\d t}f(\alpha_t) = \int \dotp{\Wgrad f(\alpha_t)(x)}{v_t(x)}\,\d\alpha_t(x) = -\int \norm{\Wgrad f(\alpha_t)(x)}^2\,\d\alpha_t(x).\)
Geodesic convexity along the geodesic
\(((1-s)\Id+sT_t)_\sharp\alpha_t\) gives
\(f(\alpha^\star)-f(\alpha_t) \geq \int \dotp{\Wgrad f(\alpha_t)(x)}{T_t(x)-x}\,\d\alpha_t(x).\)
Since \(v_t=-\Wgrad f(\alpha_t)\), this reads \(f(\alpha_t)-f(\alpha^\star) \leq \int \dotp{v_t(x)}{T_t(x)-x}\,\d\alpha_t(x).\)
\(\frac{\d}{\d t}\frac12\Wass_2^2(\alpha_t,\alpha^\star) = \int \dotp{x-T_t(x)}{v_t(x)}\,\d\alpha_t(x),\)
which proves the differential inequality. Integrating it from \(0\) to \(t\) and using the monotonicity of \(s\mapsto f(\alpha_s)\) gives
\[
t\bigl(f(\alpha_t)-f(\alpha^\star)\bigr)
\leq
\int_0^t \bigl(f(\alpha_s)-f(\alpha^\star)\bigr)\,\d s
\leq
\frac12\Wass_2^2(\alpha_0,\alpha^\star).
\]
\end{proof}
\paragraph{Convergence: the Wasserstein--PL viewpoint.}
\(|\partial f|^2(\alpha) \eqdef \int\norm{\Wgrad f(\alpha)(x)}^2\,\d\alpha(x)\)
\begin{defn}[Wasserstein--Polyak--{\L}ojasiewicz inequality]\label{def-wasserstein-pl}
Let \(f_\star\eqdef\inf_{\alpha\in\Pp_2(\RR^d)} f(\alpha)>-\infty\). The functional \(f\) satisfies the Wasserstein--PL inequality with constant \(\kappa>0\) if
\begin{equation}\label{eq-wasserstein-pl}
|\partial f|^2(\alpha)
\geq
2\kappa\bigl(f(\alpha)-f_\star\bigr)
\end{equation}
for every \(\alpha\) in the domain of \(f\). In the smooth Otto notation this reads \(\int\norm{\Wgrad f(\alpha)}^2\,\d\alpha \geq 2\kappa\bigl(f(\alpha)-f_\star\bigr).\)
\end{defn}
\begin{thm}[Wasserstein--PL convergence]\label{thm-wasserstein-pl-convergence}
Assume that \(f_\star>-\infty\), that \(f\) satisfies the Wasserstein--PL inequality~\eqref{eq-wasserstein-pl} with constant \(\kappa>0\), and that \((\alpha_t)_{t\geq0}\) is a global Wasserstein gradient flow satisfying the full energy-dissipation identity
\begin{equation}\label{eq-full-wasserstein-energy-dissipation}
\frac{\d}{\d t}f(\alpha_t)
=
-|\dot\alpha_t|^2
=
-|\partial f|^2(\alpha_t)
\qquad\text{for a.e. }t>0.
\end{equation}
Set \(E(t)\eqdef f(\alpha_t)-f_\star\). Then, for \(0\leq s\leq t\), \(E(t)\leq e^{-2\kappa(t-s)}E(s),\)
and the length of the flow segment satisfies \(\operatorname{Length}_{\Wass_2}\bigl((\alpha_r)_{r\in[s,t]}\bigr) \leq \sqrt{\frac{2}{\kappa}}\bigl(\sqrt{E(s)}-\sqrt{E(t)}\bigr).\)
If, in addition, \(f\) is lower semicontinuous for \(\Wass_2\)-convergence, then \(\alpha_t\) converges in \(\Wass_2\) to a minimizer \(\alpha_\infty\in\operatorname{Argmin} f\), and
\(\Wass_2(\alpha_t,\alpha_\infty) \leq \sqrt{\frac{2}{\kappa}}\sqrt{E(t)} \leq \sqrt{\frac{2}{\kappa}}\sqrt{E(0)}\,e^{-\kappa t}.\)
In particular, \(\operatorname{dist}_{\Wass_2}(\alpha_t,\operatorname{Argmin}f) \leq \sqrt{\frac{2}{\kappa}}\sqrt{E(0)}\,e^{-\kappa t}.\)
\end{thm}
\begin{proof}
The energy estimate is immediate from~\eqref{eq-full-wasserstein-energy-dissipation} and~\eqref{eq-wasserstein-pl}: \(E'(t) = -|\partial f|^2(\alpha_t) \leq -2\kappa E(t)\)
for almost every \(t\). Gronwall's lemma gives \(E(t)\leq e^{-2\kappa(t-s)}E(s)\).

The distance estimate uses the same two identities in a slightly different way. On the set where \(E(r)>0\), the PL inequality gives \(|\partial f|(\alpha_r)\geq \sqrt{2\kappa E(r)}.\)
Using the full dissipation identity and \(|\dot\alpha_r|=|\partial f|(\alpha_r)\),
\[
\begin{aligned}
\operatorname{Length}_{\Wass_2}\bigl((\alpha_r)_{r\in[s,t]}\bigr)
&=
\int_s^t|\dot\alpha_r|\,\d r
=
\int_s^t|\partial f|(\alpha_r)\,\d r \\
&=
\int_s^t
\frac{|\partial f|^2(\alpha_r)}{|\partial f|(\alpha_r)}\,\d r
\leq
\frac{1}{\sqrt{2\kappa}}
\int_s^t
\frac{|\partial f|^2(\alpha_r)}{\sqrt{E(r)}}\,\d r \\
&=
\frac{1}{\sqrt{2\kappa}}
\int_s^t
\frac{-E'(r)}{\sqrt{E(r)}}\,\d r
=
\sqrt{\frac{2}{\kappa}}\bigl(\sqrt{E(s)}-\sqrt{E(t)}\bigr).
\end{aligned}
\]
If \(E\) vanishes on a subinterval, then the flow is stationary there by~\eqref{eq-full-wasserstein-energy-dissipation}, so the same estimate follows by applying the argument on the part where \(E>0\). Since Wasserstein distance is bounded by metric length, the same upper bound controls \(\Wass_2(\alpha_s,\alpha_t)\).

Letting \(t\to+\infty\), the energy estimate gives \(E(t)\to0\), and the length estimate shows that the tail of \((\alpha_t)_t\) is Cauchy. Since \((\Pp_2(\RR^d),\Wass_2)\) is complete, there is a limit \(\alpha_\infty\). Lower semicontinuity yields
\(f(\alpha_\infty) \leq \liminf_{t\to+\infty}f(\alpha_t) = f_\star,\)
hence \(\alpha_\infty\in\operatorname{Argmin}f\). Sending the endpoint of the length estimate to \(+\infty\) with \(s=t\) fixed gives \(\Wass_2(\alpha_t,\alpha_\infty) \leq \sqrt{\frac{2}{\kappa}}\sqrt{E(t)}.\)
The distance-to-the-set bound follows because \(\alpha_\infty\) is one minimizer.
\end{proof}
\paragraph{Geodesic strong convexity.}
\begin{prop}[Strong geodesic convexity implies Wasserstein--PL]\label{prop-strong-geodesic-convexity-implies-pl}
Assume that \(f\) is \(\lambda\)-geodesically convex with \(\lambda>0\), that \(f_\star\) is attained at some \(\alpha^\star\), and that the Wasserstein slope is finite at \(\alpha\). Then
\(|\partial f|^2(\alpha) \geq 2\lambda\bigl(f(\alpha)-f_\star\bigr).\)
Thus positive Wasserstein strong convexity implies the Wasserstein--PL inequality with the same constant.
\end{prop}
\begin{proof}
Fix \(\alpha\) and a minimizer \(\alpha^\star\). Set
\(g=f(\alpha)-f_\star,\qquad r=\Wass_2(\alpha,\alpha^\star),\qquad a=|\partial f|(\alpha).\)
If \(r=0\), then \(\alpha=\alpha^\star\) and there is nothing to prove. Otherwise let \((\alpha_s)_{s\in[0,1]}\) be a constant-speed geodesic from \(\alpha\) to \(\alpha^\star\). By \(\lambda\)-geodesic convexity,
\(f(\alpha_s) \leq (1-s)f(\alpha)+s f_\star -\frac{\lambda}{2}s(1-s)r^2.\)
Therefore
\(f(\alpha)-f(\alpha_s) \geq s g+\frac{\lambda}{2}s(1-s)r^2.\)
Dividing by \(\Wass_2(\alpha,\alpha_s)=sr\) and letting \(s\downarrow0\) in the definition of the descending metric slope yields
\(a\geq \frac{g}{r}+\frac{\lambda}{2}r.\)
Equivalently, \(g\leq ar-\lambda r^2/2\). Maximizing the right-hand side over \(r\geq0\) gives \(g\leq a^2/(2\lambda)\), which is the desired PL inequality.
\end{proof}
\paragraph{Wasserstein Kurdyka--{\L}ojasiewicz inequalities.}
\begin{defn}[Wasserstein Kurdyka--{\L}ojasiewicz inequality]\label{def-wasserstein-kl}
Let \(f_\star=\inf f\), let \(E(\alpha)=f(\alpha)-f_\star\), and fix \(c>0\), \(\theta\in[0,1)\). We say that \(f\) satisfies a power Wasserstein--KL inequality on a set \(\mathcal U\subset \Pp_2(\RR^d)\) if
\(|\partial f|(\alpha) \geq c\,E(\alpha)^\theta, \qquad \alpha\in\mathcal U, \quad 0<E(\alpha)<+\infty.\)
Equivalently, the desingularizing function
\(\psi(s)=\frac{s^{1-\theta}}{c(1-\theta)}\)
satisfies \(\psi'(E(\alpha))|\partial f|(\alpha)\geq 1\). The Wasserstein--PL inequality \eqref{eq-wasserstein-pl} is the special case \(\theta=1/2\) with \(c=\sqrt{2\kappa}\). The regime \(\theta>1/2\) gives sublinear rates, while \(\theta=1/2\) is exactly the exponential PL regime.
\end{defn}
\begin{thm}[Sublinear convergence under Wasserstein--KL]\label{thm-wasserstein-kl-sublinear}
Let \((\alpha_t)_{t\geq0}\) be a Wasserstein gradient flow of \(f\) satisfying the full energy-dissipation identity \eqref{eq-full-wasserstein-energy-dissipation}. Assume that \(f_\star> -\infty\), that \(E(t)=f(\alpha_t)-f_\star\) is finite, and that the trajectory stays in a region where the Wasserstein--KL inequality of Definition~\ref{def-wasserstein-kl} holds with \(\theta\in(1/2,1)\). If \(E(0)>0\), then
\[
E(t)
\leq
\left(E(0)^{1-2\theta}+c^2(2\theta-1)t\right)^{-1/(2\theta-1)}.
\]
Moreover, for every \(t\geq0\), \(\operatorname{Length}_{\Wass_2}\bigl((\alpha_s)_{s\geq t}\bigr) \leq \frac{E(t)^{1-\theta}}{c(1-\theta)}.\)
If \(f\) is lower semicontinuous on \((\Pp_2(\RR^d),\Wass_2)\), then \(\alpha_t\) converges to some \(\alpha_\infty\in\operatorname{Argmin}f\) and
\[
\Wass_2(\alpha_t,\alpha_\infty)
\leq
\frac{E(t)^{1-\theta}}{c(1-\theta)}
=
O\!\left(t^{-(1-\theta)/(2\theta-1)}\right).
\]
\end{thm}
\begin{proof}
If \(E\) reaches zero at some time, then the flow is already at a global minimizer. Indeed, the integrated energy-dissipation identity forces both \(|\dot\alpha_t|\) and \(|\partial f|(\alpha_t)\) to vanish afterwards. We therefore argue on intervals where \(E>0\). For a.e. \(t\),
\(-E'(t)=|\partial f|^2(\alpha_t)=|\dot\alpha_t|^2.\)
The Wasserstein--KL inequality yields
\(E'(t) \leq -c^2E(t)^{2\theta}.\)
Since \(1-2\theta<0\), \(\frac{\d}{\d t}E(t)^{1-2\theta} = (1-2\theta)E(t)^{-2\theta}E'(t) \geq c^2(2\theta-1),\)
and integration from \(0\) to \(t\) gives the announced polynomial bound.

For the length estimate, use again \(|\dot\alpha_t|=|\partial f|(\alpha_t)\). For \(T>t\),
\[
\begin{aligned}
\int_t^T |\dot\alpha_s|\,\d s
&=
\int_t^T \frac{|\partial f|^2(\alpha_s)}{|\partial f|(\alpha_s)}\,\d s
\leq
\frac1c\int_t^T -E'(s)E(s)^{-\theta}\,\d s \\
&=
\frac{E(t)^{1-\theta}-E(T)^{1-\theta}}{c(1-\theta)}
\leq
\frac{E(t)^{1-\theta}}{c(1-\theta)}.
\end{aligned}
\]
Letting \(T\to+\infty\) shows that the tail has finite length. Hence \((\alpha_t)_t\) is Cauchy in \((\Pp_2(\RR^d),\Wass_2)\), which is complete, and converges to some \(\alpha_\infty\). Since the energy bound gives \(E(t)\to0\), lower semicontinuity gives \(f(\alpha_\infty)=f_\star\). The distance to \(\alpha_\infty\) is bounded by the remaining tail length, which proves the last estimate.
\end{proof}
\begin{prop}[Powers of PL energies are KL]\label{prop-powers-pl-are-kl}
Let \(g\geq0\) satisfy \(\inf g=0\) and the Wasserstein--PL inequality
\(|\partial g|^2(\alpha) \geq 2\kappa g(\alpha)\)
on a region where \(g>0\). Fix \(r\geq1\) and set \(f(\alpha)=g(\alpha)^r\). Whenever the metric-slope chain rule \(|\partial(g^r)|(\alpha)=r g(\alpha)^{r-1}|\partial g|(\alpha)\) holds, \(f\) satisfies the Wasserstein--KL inequality
\(|\partial f|(\alpha) \geq r\sqrt{2\kappa}\, f(\alpha)^{1-1/(2r)}.\)
Thus \(f\) has KL exponent \(\theta=1-1/(2r)\). For \(r=1\) this is PL; for \(r>1\), Theorem~\ref{thm-wasserstein-kl-sublinear} gives \(f(\alpha_t)=O(t^{-r/(r-1)})\) along the \(f\)-gradient flow, under its hypotheses.
\end{prop}
\begin{proof}
By the chain rule and the PL inequality for \(g\), \(|\partial f| = rg^{r-1}|\partial g| \geq r\sqrt{2\kappa}\,g^{r-1/2} = r\sqrt{2\kappa}\,f^{1-1/(2r)}.\)
The exponent and rate follow by substituting \(\theta=1-1/(2r)\) into Theorem~\ref{thm-wasserstein-kl-sublinear}.
\end{proof}
\paragraph{Convexity and curvature.}
The same language is not restricted to subsets of $\RR^d$. If $(\X,\dist,\mathfrak m)$ is a geodesic metric-measure space, $\Wass_2$ geodesics can be defined by transporting each pair of endpoints along metric geodesics, or more intrinsically by dynamical optimal plans on path space, as discussed in Section~\ref{sec-kantorovich-plan-interpolation}.

\[
\mathrm{Ent}_{\mathfrak m}(\alpha)
\eqdef
\begin{cases}
\displaystyle \int_\X \rho\log\rho\,\d\mathfrak m,
& \text{if } \alpha=\rho\,\mathfrak m,\\
+\infty,
& \text{otherwise.}
\end{cases}
\]
\begin{thm}[Ricci curvature and entropy convexity]\label{thm-ricci-entropy-convexity}
Let $(M,g)$ be a smooth compact connected Riemannian manifold without boundary, and let $\mathfrak m=\mathrm{vol}_g$. For $\lambda\in\RR$, the lower Ricci bound $\mathrm{Ric}_g\geq\lambda g$ holds if and only if $\mathrm{Ent}_{\mathfrak m}$ is $\lambda$-geodesically convex on $(\Pp_2(M),\Wass_2)$.
\end{thm}
\subsection{Functional Inequalities via Optimal Transport}
Optimal transport proves inequalities by turning comparison into geometry. One transports a density by a monotone or Brenier map, interpolates along the resulting rays, and converts concavity of a Jacobian determinant or convexity of an energy into an integral estimate.

\paragraph{Brunn--Minkowski and isoperimetry.}
\begin{prop}[Brunn--Minkowski and Euclidean isoperimetry]\label{prop-ot-brunn-minkowski-isoperimetric}
Let $A,B\subset\RR^d$ be bounded Borel sets with positive Lebesgue measure. For $t\in[0,1]$, set
\((1-t)A+tB\eqdef\{(1-t)x+ty\;:\;x\in A,\ y\in B\}.\)
Then
\(\mathrm{Vol}((1-t)A+tB)^{1/d} \geq (1-t)\mathrm{Vol}(A)^{1/d} + t\,\mathrm{Vol}(B)^{1/d}.\)
\(\mathrm{Per}(A) \geq d\,\omega_d^{1/d}\mathrm{Vol}(A)^{(d-1)/d}.\)
The constant is sharp, with equality for Euclidean balls.
\end{prop}
\begin{proof}
We first give the argument for regular sets, the general Borel case following by approximation from inside and outside. Let $\alpha$ and $\beta$ be the uniform probability measures on $A$ and $B$, and let $T=\nabla\phi$ be the Brenier map with $T_\sharp\alpha=\beta$. The interpolating map $T_t=(1-t)\Id+tT$ sends $\alpha$ to a measure $\alpha_t$ supported in $(1-t)A+tB$. At points where $T$ is differentiable,
\(\det DT_t = \det((1-t)\Id+tDT).\)
Since $DT$ is symmetric positive semidefinite and $\det^{1/d}$ is concave on the positive semidefinite cone, \(\det(DT_t)^{1/d} \geq (1-t)+t\det(DT)^{1/d}.\)
The change-of-variables identity for the uniform measures gives $\det(DT)=\mathrm{Vol}(B)/\mathrm{Vol}(A)$ almost everywhere. If $\rho_t$ denotes the density of $\alpha_t$, a second application of the same formula gives
\(\rho_t(T_t(x))^{-1/d} = \mathrm{Vol}(A)^{1/d}\det(DT_t(x))^{1/d} \geq c_t,\)
where $c_t=(1-t)\mathrm{Vol}(A)^{1/d}+t\mathrm{Vol}(B)^{1/d}$. Thus $\rho_t\leq c_t^{-d}$ almost everywhere. Since $\alpha_t$ has total mass one and is supported in $(1-t)A+tB$, this gives $\mathrm{Vol}((1-t)A+tB)\geq c_t^d$, which is the displayed Brunn--Minkowski inequality.

For isoperimetry, use the homogeneity of Brunn--Minkowski and apply it to $A$ and $\epsilon B_1$. Writing $A+\epsilon B_1$ for the parallel set gives
\(\mathrm{Vol}(A+\epsilon B_1)^{1/d} \geq \mathrm{Vol}(A)^{1/d}+\epsilon\omega_d^{1/d}.\)
Using $\mathrm{Vol}(A+\epsilon B_1)=\mathrm{Vol}(A)+\epsilon\,\mathrm{Per}(A)+o(\epsilon)$ and differentiating at $\epsilon=0$ yields the claimed perimeter bound. Balls attain equality, so the constant is optimal.
\end{proof}
\paragraph{Pr\'ekopa--Leindler.}
\begin{prop}[Pr\'ekopa--Leindler inequality]\label{prop-ot-prekopa-leindler}
Let $u,v,w:\RR^d\to[0,+\infty)$ be integrable, and let $t\in[0,1]$. Assume that
\(w((1-t)x+ty)\geq u(x)^{1-t}v(y)^t \qquad \text{for all }x,y\in\RR^d.\)
Then
\[
\int_{\RR^d} w
\geq
\left(\int_{\RR^d}u\right)^{1-t}
\left(\int_{\RR^d}v\right)^t.
\]
\end{prop}
\begin{proof}
The endpoint cases are immediate, so fix $t\in(0,1)$. Set $U=\int u$ and $V=\int v$. If $U=0$ or $V=0$, the conclusion is trivial; hence assume $U,V>0$. Let $\alpha=(u/U)\d x$ and $\beta=(v/V)\d x$, and let $T$ be the Brenier map from $\alpha$ to $\beta$. Put $T_t=(1-t)\Id+tT$. On a smooth positive regularization, $T_t$ is the gradient of a strictly convex function and is therefore one-to-one almost everywhere. The area formula gives
\(\int w \geq \int w(T_t(x))\det(DT_t(x))\d x.\)
For each eigenvalue $s\geq0$ of $DT$, weighted arithmetic--geometric mean gives $(1-t)+ts\geq s^t$. Hence
\(\det((1-t)\Id+tDT)\geq\det(DT)^t.\)
\(\int w \geq \int u(x)^{1-t}v(T(x))^t\det(DT(x))^t\d x.\)
Since $T_\sharp\alpha=\beta$, the Monge--Amp\`ere identity gives \(\det(DT(x))=\frac{u(x)V}{v(T(x))U}\)
almost everywhere. Substitution yields
\[
\int w
\geq
\left(\frac{V}{U}\right)^t\int u(x)\d x
=
U^{1-t}V^t.
\]
Approximation removes the smoothness and positivity assumptions.
\end{proof}
\paragraph{Logarithmic Sobolev inequalities.}
\begin{equation}\label{eq-relative-fisher-information}
\mathcal I(\alpha|\beta)
\eqdef
\int \norm{\nabla\log h}^2\,\d\alpha
=
\int \frac{\norm{\nabla h}^2}{h}\,\d\beta,
\end{equation}
\begin{equation}\label{eq-general-log-sobolev}
\KL(\alpha|\beta)
\leq
\frac{1}{2\lambda}\mathcal I(\alpha|\beta)
\qquad
\text{for all probability measures }\alpha.
\end{equation}
\begin{prop}[Curvature implies logarithmic Sobolev]\label{prop-curvature-log-sobolev}
Let \(\beta=Z^{-1}e^{-V}\d x\) be a smooth probability measure on \(\RR^d\). If \(\nabla^2V\succeq\lambda\Id\) for some \(\lambda>0\), then \(\beta\) satisfies \eqref{eq-general-log-sobolev} with constant \(\lambda\).
\end{prop}
\begin{proof}
Apply the HWI inequality of Proposition~\ref{prop-hwi-inequality}. Under the curvature assumption, for every absolutely continuous \(\alpha\),
\(\KL(\alpha|\beta) \leq \Wass_2(\alpha,\beta)\sqrt{\mathcal I(\alpha|\beta)} - \frac{\lambda}{2}\Wass_2^2(\alpha,\beta).\)
For fixed \(s=\sqrt{\mathcal I(\alpha|\beta)}\), maximize the right-hand side over \(r=\Wass_2(\alpha,\beta)\geq0\). The elementary inequality \(rs-\lambda r^2/2\leq s^2/(2\lambda)\) gives \eqref{eq-general-log-sobolev}. Approximation yields the standard nonsmooth statement. This is the Otto--Villani transport route from the Bakry--Emery curvature bound to logarithmic Sobolev; the original \(\Gamma_2\) calculus gives the same criterion directly.
\end{proof}
\begin{prop}[Logarithmic Sobolev inequality implies entropy decay]\label{prop-lsi-entropy-decay}
Let \(\beta=e^{-V}\d x/Z\) be a smooth probability measure on \(\RR^d\) satisfying the logarithmic Sobolev inequality \eqref{eq-general-log-sobolev} with constant \(\lambda>0\). Let \(\alpha_t\) solve the Wasserstein gradient flow of \(f(\alpha)=\KL(\alpha|\beta)\), namely
\(\partial_t\rho_t = \nabla\cdot(\rho_t\nabla V)+\Delta\rho_t, \qquad \alpha_t=\rho_t\d x.\)
If \(\KL(\alpha_0|\beta)<+\infty\), then
\(\KL(\alpha_t|\beta) \leq e^{-2\lambda t}\KL(\alpha_0|\beta).\)
Moreover, if \(\alpha_0,\beta\in\Pp_2(\RR^d)\), then
\(\Wass_2(\alpha_t,\beta) \leq \sqrt{\frac{2}{\lambda}\KL(\alpha_0|\beta)}\,e^{-\lambda t}.\)
\end{prop}
\begin{proof}
The first variation of \(f\) is \(\log(\rho/\rho_\beta)\), up to an additive constant. Along the smooth Wasserstein gradient flow,
\(\frac{\d}{\d t}\KL(\alpha_t|\beta) = -\int \norm{\nabla\log\frac{\d\alpha_t}{\d\beta}}^2\,\d\alpha_t = -\mathcal I(\alpha_t|\beta).\)
The logarithmic Sobolev inequality bounds the Fisher information below by \(2\lambda\KL(\alpha_t|\beta)\). Hence
\(\frac{\d}{\d t}\KL(\alpha_t|\beta) \leq -2\lambda\KL(\alpha_t|\beta),\)
and Gronwall's lemma gives the entropy estimate. The Wasserstein estimate follows from the length estimate in the Wasserstein--PL theorem: apply Theorem~\ref{thm-wasserstein-pl-convergence} to \(f=\KL(\cdot|\beta)\), whose unique minimizer is \(\beta\), and use the logarithmic Sobolev inequality as the corresponding Wasserstein--PL inequality.
\end{proof}
\paragraph{Gaussian logarithmic Sobolev inequality.}
\begin{prop}[Gaussian logarithmic Sobolev inequality]\label{prop-gaussian-log-sobolev-ot}
Let \(\gamma_d\) be the standard Gaussian measure on \(\RR^d\). If \(\rho>0\) is smooth, has sufficient decay, and \(\int \rho\,\d\gamma_d=1\), then
\(\mathrm{Ent}_{\gamma_d}(\rho) \eqdef \int \rho\log \rho\,\d\gamma_d \leq \frac12\int \frac{\norm{\nabla \rho}^2}{\rho}\,\d\gamma_d.\)
\end{prop}
\begin{proof}
Let \(\alpha=\rho\gamma_d\), and let \(T=\nabla\phi\) be the Brenier map with \(T_\sharp\alpha=\gamma_d\). Write \(T(x)=x+\xi(x)\). The Monge--Amp\`ere equation reads \(\rho(x)e^{-\norm{x}^2/2} = e^{-\norm{T(x)}^2/2}\det(DT(x)).\)
Thus
\(\log \rho(x) = -\dotp{x}{\xi(x)} -\frac12\norm{\xi(x)}^2 + \log\det(\Id+D\xi(x)).\)
Since \(DT\) is positive semidefinite, \(\log\det(\Id+D\xi)\leq\tr(D\xi)=\diverg\xi\). Multiplying by \(\rho\) and integrating with respect to \(\gamma_d\) gives
\(\mathrm{Ent}_{\gamma_d}(\rho) \leq \int \rho(\diverg\xi-\dotp{x}{\xi})\,\d\gamma_d -\frac12\int \rho\norm{\xi}^2\,\d\gamma_d.\)
Gaussian integration by parts gives \(\int \rho(\diverg\xi-\dotp{x}{\xi})\,\d\gamma_d = -\int\dotp{\nabla \rho}{\xi}\,\d\gamma_d.\)
Completing the square pointwise, \(-\dotp{\nabla \rho}{\xi}-\frac12 \rho\norm{\xi}^2 \leq \frac12\frac{\norm{\nabla \rho}^2}{\rho}.\)
This proves the inequality. The general Sobolev-class statement follows by approximation.
\end{proof}
\[
\rho_a(x)=\exp\!\left(\dotp{a}{x}-\frac12\norm{a}^2\right),
\qquad a\in\RR^d,
\]
\(\int \rho_a\log \rho_a\,\d\gamma_d = \frac12\norm{a}^2, \qquad \frac12\int\frac{\norm{\nabla \rho_a}^2}{\rho_a}\,\d\gamma_d = \frac12\norm{a}^2.\)
Thus equality is attained along translations, and the sharp logarithmic Sobolev constant of the standard Gaussian is \(\lambda=1\). Proposition~\ref{prop-lsi-entropy-decay} then gives the classical exponential convergence of the Ornstein--Uhlenbeck flow.

\paragraph{Poincar\'e as a linearized Wasserstein--PL inequality.}
\(\beta_\epsilon=(1+\epsilon\xi)\beta, \qquad \beta_0=\beta.\)
\[
\mathsf H^f_\beta(\xi)
\eqdef
\lim_{\epsilon\to0}
\frac{2(f(\beta_\epsilon)-f(\beta))}{\epsilon^2},
\qquad
\mathsf D^f_\beta(\xi)
\eqdef
\lim_{\epsilon\to0}
\frac{|\partial f|^2(\beta_\epsilon)}{\epsilon^2},
\]
\begin{prop}[Linearizing a Wasserstein--PL inequality]\label{prop-linearized-pl-functional-inequality}
Let \(\beta\) be a minimizer of an energy \(f\), and assume that \(f\) satisfies \(|\partial f|^2(\alpha) \geq 2\kappa\bigl(f(\alpha)-f(\beta)\bigr)\)
in a neighborhood of \(\beta\). Let \(\xi\) be a smooth bounded perturbation with \(\int\xi\,\d\beta=0\), and set \(\beta_\epsilon=(1+\epsilon\xi)\beta\). If \(\mathsf H^f_\beta(\xi)\) and \(\mathsf D^f_\beta(\xi)\) exist, then
\(\kappa\,\mathsf H^f_\beta(\xi) \leq \mathsf D^f_\beta(\xi).\)
\end{prop}
\begin{proof}
Apply the PL inequality to \(\beta_\epsilon\): \(|\partial f|^2(\beta_\epsilon) \geq 2\kappa\bigl(f(\beta_\epsilon)-f(\beta)\bigr).\)
Divide by \(\epsilon^2\) and let \(\epsilon\to0\). The definitions of \(\mathsf H^f_\beta\) and \(\mathsf D^f_\beta\) give the claim.
\end{proof}
For \(f(\alpha)=\KL(\alpha|\beta)\), Proposition~\ref{prop-linearized-pl-functional-inequality} says precisely that logarithmic Sobolev linearizes to Poincar\'e. Indeed, for \(h_\epsilon=1+\epsilon\xi\),

\[
\KL(h_\epsilon\beta|\beta)
=
\frac{\epsilon^2}{2}\int\xi^2\,\d\beta+o(\epsilon^2),
\qquad
\mathcal I(h_\epsilon\beta|\beta)
=
\epsilon^2\int\norm{\nabla\xi}^2\,\d\beta+o(\epsilon^2).
\]
\begin{prop}[Poincar\'e from logarithmic Sobolev]\label{prop-poincare-linearized-wasserstein-pl}
Let \(\beta\) be a smooth probability measure on \(\RR^d\). Assume that \(\beta\) satisfies the logarithmic Sobolev inequality with constant \(\lambda>0\),
\[
\KL(\alpha|\beta)
\leq
\frac{1}{2\lambda}\,\mathcal I(\alpha|\beta),
\qquad
\mathcal I(\alpha|\beta)
\eqdef
\int\norm{\nabla\log\frac{\d\alpha}{\d\beta}}^2\,\d\alpha.
\]
Then, for every smooth \(\xi\) with \(\int\xi\,\d\beta=0\), \(\lambda\int\xi^2\,\d\beta \leq \int\norm{\nabla\xi}^2\,\d\beta.\)
In particular, if \(\d\beta=Z^{-1}e^{-V}\d x\) and \(\nabla^2 V\succeq \lambda\Id\), so that \(\beta\) is \(\lambda\)-strongly log-concave, then the conclusion holds by the Bakry--Emery criterion.
\end{prop}
\begin{proof}
For \(h_\epsilon=1+\epsilon\xi\), with \(|\epsilon|\) small enough, one has \(\beta_\epsilon=h_\epsilon\beta\) and
\(\KL(\beta_\epsilon|\beta) = \int h_\epsilon\log h_\epsilon\,\d\beta = \frac{\epsilon^2}{2}\int\xi^2\,\d\beta+o(\epsilon^2),\)
because \(\int\xi\,\d\beta=0\). The Fisher information expands as
\[
\mathcal I(\beta_\epsilon|\beta)
=
\int\norm{\nabla\log h_\epsilon}^2 h_\epsilon\,\d\beta
=
\epsilon^2\int\norm{\nabla\xi}^2\,\d\beta+o(\epsilon^2).
\]
Inserting these expansions in the logarithmic Sobolev inequality and letting \(\epsilon\to0\) gives the Poincar\'e inequality. The last assertion is the usual Bakry--Emery implication from \(\lambda\)-convexity of \(V\) to the logarithmic Sobolev inequality.
\end{proof}
\(L_\beta\xi \eqdef \Delta\xi-\dotp{\nabla V}{\nabla\xi} = \rho_\beta^{-1}\nabla\cdot(\rho_\beta\nabla\xi).\)
\(\int\norm{\nabla\xi}^2\,\d\beta = \langle \xi,-L_\beta\xi\rangle_{L^2(\beta)}.\)
\[
\norm{\dot h}_{-1,\beta}^2
\eqdef
\inf_{v}
\left\{
\int\norm v^2\,\d\beta
\;:\;
\dot h=-\rho_\beta^{-1}\nabla\cdot(\rho_\beta v)
\right\}.
\]
\paragraph{Talagrand's transport-entropy inequality.}
The distance estimate in Theorem~\ref{thm-wasserstein-pl-convergence} already reveals the general principle: when \(\KL(\cdot|\beta)\) satisfies logarithmic Sobolev with constant \(\lambda\), the entropy flow yields the transport-entropy bound \(\Wass_2^2(\alpha,\beta)\leq2\KL(\alpha|\beta)/\lambda\).

\begin{prop}[Gaussian \(T_2\) inequality]\label{prop-gaussian-talagrand-t2}
For every $\alpha\in\Pp_2(\RR^d)$, \(\frac12\Wass_2^2(\alpha,\gamma_d) \leq \KL(\alpha|\gamma_d),\)
with the convention that the right-hand side is \(+\infty\) when \(\alpha\not\ll\gamma_d\).
\end{prop}
\begin{proof}
If \(\KL(\alpha|\gamma_d)=+\infty\), the claim is immediate. Assume first that \(\alpha=\rho\gamma_d\) with \(\rho\) smooth and positive. Let $T=\nabla\phi$ be the Brenier map with $T_\sharp\gamma_d=\alpha$, and write $T(x)=x+\xi(x)$. The change of variables gives
\(\rho(T(x))e^{-\norm{T(x)}^2/2}\det(DT(x)) = e^{-\norm{x}^2/2}.\)
Hence
\(\log \rho(T(x)) = \dotp{x}{\xi(x)} + \frac12\norm{\xi(x)}^2 - \log\det(\Id+D\xi(x)).\)
Using again $\log\det(\Id+D\xi)\leq\diverg\xi$ and integrating with respect to $\gamma_d$, \(\KL(\alpha|\gamma_d) \geq \frac12\int\norm{\xi}^2\,\d\gamma_d + \int(\dotp{x}{\xi}-\diverg\xi)\,\d\gamma_d.\)
The last integral vanishes by Gaussian integration by parts. Since $T$ is optimal, \(\int\norm{\xi}^2\,\d\gamma_d=\Wass_2^2(\gamma_d,\alpha),\)
which proves the claim. Approximation gives the general case. Equality holds for translations \(\alpha=\Gaussian(a,\Id)\), so the constant is sharp.
\end{proof}
\paragraph{The HWI bridge.}
\begin{prop}[Otto--Villani HWI inequality]\label{prop-hwi-inequality}
Let $\beta=\rho_\beta\d x=e^{-V}\d x/Z$ be a smooth probability measure in $\Pp_2(\RR^d)$, and assume that $\nabla^2 V\succeq \lambda \Id$ for some $\lambda\in\RR$. For $\alpha\in\Pp_2(\RR^d)$, define the relative Fisher information by
\(\mathcal I(\alpha|\beta) \eqdef \int \norm{\nabla\log\frac{\rho_\alpha}{\rho_\beta}}^2\,\d\alpha,\)
when \(\alpha=\rho_\alpha\d x\) and the expression is finite, and set \(\mathcal I(\alpha|\beta)=+\infty\) otherwise. Then
\(\KL(\alpha|\beta) \leq \Wass_2(\alpha,\beta)\sqrt{\mathcal I(\alpha|\beta)} -\frac{\lambda}{2}\Wass_2^2(\alpha,\beta).\)
\end{prop}
\begin{proof}
If $\mathcal I(\alpha|\beta)=+\infty$, there is nothing to prove. Assume first that $\rho_\alpha$ is smooth and positive. Let $T$ be the Brenier map from $\alpha$ to $\beta$, set $v=T-\Id$, and let $\alpha_t=((1-t)\Id+tT)_\sharp\alpha$. By Theorem~\ref{thm-mccann-internal-energy} and Proposition~\ref{prop-basic-geodesic-convexity}, $f(\eta)=\KL(\eta|\beta)$ is $\lambda$-geodesically convex. Thus the one-dimensional function $e(t)=f(\alpha_t)$ satisfies
\(e(1)\geq e(0)+e'(0)+\frac{\lambda}{2}\Wass_2^2(\alpha,\beta).\)
Since $e(1)=f(\beta)=0$, it remains to compute the initial derivative. The continuity equation for the geodesic gives $\partial_t\alpha_t+\nabla\cdot(\alpha_t v_t)=0$ and $v_0=v$. Hence
\(e'(0) = \int \dotp{\nabla\log\frac{\rho_\alpha}{\rho_\beta}}{v}\,\d\alpha.\)
Therefore
\(\KL(\alpha|\beta) \leq -\int \dotp{\nabla\log\frac{\rho_\alpha}{\rho_\beta}}{T-\Id}\,\d\alpha -\frac{\lambda}{2}\Wass_2^2(\alpha,\beta).\)
Cauchy--Schwarz and $\int\norm{T-\Id}^2\d\alpha=\Wass_2^2(\alpha,\beta)$ give the announced bound. The general case follows by standard approximation and lower semicontinuity of entropy, Fisher information, and $\Wass_2$.
\end{proof}
When \(\lambda>0\), the HWI inequality is a bridge from curvature to PL. It is not itself a PL inequality because it still contains the distance term \(\Wass_2(\alpha,\beta)\). Optimizing the right-hand side, or simply using Young's inequality \(rs-\lambda r^2/2\leq s^2/(2\lambda)\), gives the logarithmic Sobolev estimate

\(\KL(\alpha|\beta) \leq \frac{1}{2\lambda}\mathcal I(\alpha|\beta).\)
\subsection{Training Two-Layer MLPs as Wasserstein Flows}
\label{sec-wasserstein-flows-mlp}
\paragraph{Wasserstein training of two-layer MLPs.}
\(x=(u,v)\in\RR^d\times\RR^{d'},\)
where $u$ is the inner weight and $v$ is the outer vector weight. For a scalar nonlinearity $\sigma$, define the vector-valued feature

\(\psi(x,z)=v\,\sigma(\dotp{u}{z})\in\RR^{d'}.\)
\(G_X(z)=\frac1n\sum_{i=1}^n\psi(x_i,z), \qquad G_\alpha(z)=\int\psi(x,z)\d\alpha(x), \qquad \alpha=\frac1n\sum_i\delta_{x_i}.\)
Let $\zeta$ be a probability distribution on data-label pairs $(z,y)\in\RR^d\times\RR^{d'}$. The population risk is

\(f(\alpha)=\int \ell(G_\alpha(z),y)\d\zeta(z,y),\)
\[
\dot x_i=-n\nabla_{x_i}F(X),
\qquad
F(X)=f\!\left(\frac1n\sum_i\delta_{x_i}\right),
\]
which is the gradient flow of $F(X)=f(\alpha_X)$ for this Wasserstein particle metric, equivalently Euclidean gradient descent with the time scale multiplied by $n$. It gives a particle discretization of~\eqref{eq:wassflow-pde}.

Assume that $\ell$ is differentiable in its first variable. The first variation is

\begin{equation}\label{eq-mlp-first-variation-general}
\delta f(\alpha)(x)
=
\int
\dotp{\nabla_1\ell(G_\alpha(z),y)}{\psi(x,z)}
\d\zeta(z,y),
\end{equation}
\(\Wgrad f(\alpha)(x) = \nabla_x\delta f(\alpha)(x) = \int [D_x\psi(x,z)]^\top\nabla_1\ell(G_\alpha(z),y) \d\zeta(z,y).\)
For the squared Euclidean loss $\ell(s,y)=\frac12\norm{s-y}^2$, the energy is the sum of a quadratic interaction and a linear potential:

\begin{equation}\label{eq-mlp-square-loss-quadratic-linear}
f(\alpha)
=
\frac12\iint k(x,x')\d\alpha(x)\d\alpha(x')
+
\int g(x)\d\alpha(x)
+\frac12\int\norm{y}^2\d\zeta(z,y),
\end{equation}
\begin{equation}\label{eq-mlp-kernel-linear-potential}
k(x,x')=\int\dotp{\psi(x,z)}{\psi(x',z)}\d\zeta(z,y),
\qquad
g(x)=-\int\dotp{y}{\psi(x,z)}\d\zeta(z,y).
\end{equation}
\(\delta f(\alpha)(x)=\int k(x,x')\d\alpha(x')+g(x), \qquad \Wgrad f(\alpha)(x)=\int\nabla_x k(x,x')\d\alpha(x')+\nabla_xg(x).\)
\paragraph{Classical convexity and stationarity.}
\(f(\alpha) = \frac12\iint k(x,x')\d\alpha(x)\d\alpha(x') + \int V(x)\d\alpha(x) +C,\)
\(Q((1-s)\alpha+s\beta)\leq (1-s)Q(\alpha)+sQ(\beta), \qquad Q(\alpha)=\frac12\iint k\d\alpha\d\alpha.\)
\begin{prop}[Affine convexity and stationary positive densities]\label{prop-classical-convex-stationary}
Let $\Omega\subset\RR^d$ be connected, let $f=Q+\int V\,\d\alpha+C$ be as above on $\Pp(\Omega)$, and assume that $Q$ is classically convex. Suppose that a Wasserstein gradient flow $\alpha_t$ of $f$ converges to $\alpha_\infty=\rho_\infty\d x$, where $\rho_\infty>0$ almost everywhere on $\Omega$, and that $\inf_{t\geq0}f(\alpha_t)>-\infty$. Assume that $\delta f(\alpha_\infty)\in C^1(\Omega)$ and that the squared Wasserstein slope is lower semicontinuous along the convergence, in the sense that
\[
\int_\Omega\norm{\nabla\delta f(\alpha_\infty)}^2\d\alpha_\infty
\leq
\liminf_{n\to\infty}
\int_\Omega\norm{\nabla\delta f(\alpha_{t_n})}^2\d\alpha_{t_n}
\]
whenever $t_n\to\infty$. Then $\alpha_\infty$ is a global minimizer of $f$.
\end{prop}
\begin{proof}
The energy-dissipation identity and the lower bound on $f$ along the convergent trajectory imply
\(\int_0^\infty \int_\Omega\norm{\nabla\delta f(\alpha_t)}^2\d\alpha_t\d t <+\infty.\)
Hence there exists $t_n\to\infty$ for which the inner integral tends to zero. The assumed lower semicontinuity gives \(\int_\Omega \norm{\nabla \delta f(\alpha_\infty)}^2\d\alpha_\infty=0.\)
Because $\rho_\infty>0$ almost everywhere and $\nabla\delta f(\alpha_\infty)$ is continuous, this gradient vanishes throughout $\Omega$. Connectedness then implies that $\delta f(\alpha_\infty)$ is constant.
\(\int_\Omega\delta f(\alpha_\infty)\d(\beta-\alpha_\infty)=0.\)
Classical convexity of $f$ in the affine variable $\alpha$ gives the subgradient inequality
\(f(\beta)\geq f(\alpha_\infty) + \int \delta f(\alpha_\infty)\d(\beta-\alpha_\infty) = f(\alpha_\infty).\)
Thus no competitor has smaller energy. Strict positivity is essential to this argument: without it, stationarity controls the first variation only on the support explored by the limiting measure. For square-loss two-layer mean-field models, \eqref{eq-mlp-square-loss-quadratic-linear} is exactly of this quadratic-plus-linear form, and positive semidefiniteness of the induced kernel $k$ is the classical convexity assumption.
\end{proof}
\begin{prop}[Formal global optimality for two-homogeneous mean-field flows]\label{prop-formal-chizat-bach}
Assume that the feature is positively two-homogeneous in the neuron variable, \(\psi(\lambda x,z)=\lambda^2\psi(x,z) \qquad(\lambda>0),\)
and that $f(\alpha)=J(G_\alpha)$ with $J$ convex and differentiable as a functional of the predictor. Let $\alpha$ be a smooth stationary point of the Wasserstein flow, so that $\nabla_x\delta f(\alpha)(x)=0$ on $\supp(\alpha)$. Assume also full directional support: for every unit direction $\omega$, the support of $\alpha$ intersects the ray $\{\lambda\omega:\lambda>0\}$. Then $\alpha$ is a global minimizer of $f$ over the mean-field model class.
\end{prop}
\begin{proof}
Choose the first-variation representative
\[
h_\alpha(x)=\delta f(\alpha)(x)
=
\left\langle \nabla J(G_\alpha),\psi(x,\cdot)\right\rangle_\rho.
\]
Additive constants in first variations do not affect the Wasserstein gradient or first-order inequalities between probability measures; this normalization is useful because it inherits the homogeneity of $\psi$. By two-homogeneity of $\psi$, one has
\(h_\alpha(\lambda x)=\lambda^2h_\alpha(x), \qquad \lambda>0.\)
Taking \(x=0\) and any \(\lambda\neq1\) in this identity gives \(h_\alpha(0)=0\). We first record that stationarity forces \(h_\alpha\) to vanish on \(\supp(\alpha)\). Indeed, if \(x=r\omega\in\supp(\alpha)\) with \(r>0\) and \(\norm{\omega}=1\), then
\(0=\dotp{\nabla h_\alpha(r\omega)}{\omega} =\frac{\d}{\d s}h_\alpha(s\omega)\bigg|_{s=r} =2r h_\alpha(\omega),\)
so \(h_\alpha(\omega)=0\), and hence \(h_\alpha(x)=r^2h_\alpha(\omega)=0\). The case \(x=0\) has already been covered, so \(h_\alpha=0\) on \(\supp(\alpha)\) and \(\int h_\alpha\,\d\alpha=0\).

Suppose that a competitor \(\beta\) satisfies \(f(\beta)<f(\alpha)\). Since
\(G_\beta-G_\alpha = \int \psi(x,\cdot)\d(\beta-\alpha)(x),\)
convexity of \(J\) gives \(f(\beta)-f(\alpha) \geq \int h_\alpha(x)\d(\beta-\alpha)(x).\)
The left-hand side is negative, while \(\int h_\alpha\,\d\alpha=0\), so \(\int h_\alpha\,\d\beta<0\). Since \(h_\alpha(0)=0\), this strict negativity must occur away from the origin: there exists a nonzero point \(x=r\omega\) with \(r>0\), \(\norm{\omega}=1\), and \(h_\alpha(x)<0\). Equivalently \(h_\alpha(\omega)=h_\alpha(x)/r^2<0\). By full directional support, there is \(r_\omega>0\) such that \(r_\omega\omega\in\supp(\alpha)\). At this support point,
\(\dotp{\nabla h_\alpha(r_\omega\omega)}{\omega} = 2r_\omega h_\alpha(\omega) <0,\)
which contradicts the stationarity condition \(\nabla h_\alpha=0\) on \(\supp(\alpha)\). Thus no competitor has smaller risk.

The rigorous Chizat--Bach theorem replaces the full directional support assumption by propagation and overparameterization hypotheses ensuring that any negative first-variation direction is represented by the evolving support. The same radial-derivative contradiction, applied after this support-propagation step, then rules out non-optimal stationary limits.
\end{proof}
\subsection{Generalized Dynamic Wasserstein Flows}
\label{sec-generalized-dynamic-wasserstein-flows}
\paragraph{Generalized Wasserstein flows.}
\begin{equation}\label{eq-generalized-jko-step}
\alpha_{t+\tau}
\in
\uargmin{\alpha}
\frac{1}{2\tau}d^2(\alpha_t,\alpha)+f(\alpha).
\end{equation}
\begin{equation}\label{eq-generalized-flow-continuity-motivation}
\partial_t\alpha_t+\diverg(\alpha_t v_{\alpha_t})=0,
\end{equation}
The metric ingredient needed below is the squared-speed action $\mathbb A(\alpha,w)$ introduced in Section~\ref{sec-generalized-dynamic-wasserstein-distances}. Its length distance is already defined in~\eqref{eq-generalized-action-length-distance}; the present section only uses the local cost in the infinitesimal variational step defined below.

\[
\mathbb A_p(\alpha,w)=\left(\int\norm w^p\d\alpha\right)^{2/p}.
\]
\(A_2(a,w)=a\norm w^2, \qquad J_2(a,m)=\frac{\norm m^2}{a}, \qquad \mathsf D_{\mathbb A}=\Wass_2.\)
\begin{defn}[Penalized minimization oracle]\label{def-generalized-action-steepest-direction}
Let \(\mathbb A(\alpha,\cdot)\) be the local squared-speed action at \(\alpha\). For a cotangent field \(u\) for which the pairing below is finite, usually \(u=\Wgrad f(\alpha)\), the penalized minimization oracle associated with this action is
\begin{equation}\label{eq-local-action-steepest-descent}
\operatorname{PMO}_{\mathbb A,\alpha}(u)
\in
\uargmin{w}
\left\{
\int \dotp{u(x)}{w(x)}\d\alpha(x)
+
\frac12\mathbb A(\alpha,w)
\right\},
\end{equation}
where the minimization is over admissible finite-action velocity representatives for the chosen geometry. In Hilbertian examples \(u\in L^2(\alpha;\RR^d)\); for \(\Wass_p\), the natural dual space is \(L^q(\alpha;\RR^d)\), with \(q=p/(p-1)\).
\end{defn}
For an energy $f$, the formal small-step expansion of~\eqref{eq-generalized-jko-step}, when the distance is generated by the squared-speed action \(\mathbb A\) and the JKO displacement admits an admissible finite-action tangent representative, selects the PMO direction

\begin{equation}\label{eq-generalized-gradient-from-action}
v_{\alpha}
=
\operatorname{PMO}_{\mathbb A,\alpha}(\Wgrad f(\alpha)).
\end{equation}
\begin{equation}\label{eq-local-action-gradient-flow-pde}
\partial_t\alpha_t
+
\diverg\!\left(
\alpha_t\,\operatorname{PMO}_{\mathbb A,\alpha_t}(\Wgrad f(\alpha_t))
\right)=0.
\end{equation}
\[
\frac{\d}{\d t}f(\alpha_t)
=
\left\langle \Wgrad f(\alpha_t),v_t\right\rangle_{\alpha_t}
=
-\mathbb A(\alpha_t,v_t),
\qquad
v_t=\operatorname{PMO}_{\mathbb A,\alpha_t}(\Wgrad f(\alpha_t)).
\]
For \(\Wass_2\), \(\mathbb A(\alpha,w)=\int\norm w^2\d\alpha\) and \(\operatorname{PMO}_{\mathbb A,\alpha}(u)=-u\). Hence~\eqref{eq-local-action-gradient-flow-pde} reduces to the standard Wasserstein gradient-flow equation

\[
\partial_t\alpha_t
=
\diverg\!\left(\alpha_t\Wgrad f(\alpha_t)\right),
\]
In the restricted Riemannian case of~\eqref{eq-general-quadratic-tangent-action}, where

\[
\mathbb A(\alpha,w)=\left\langle Q_\alpha w,w\right\rangle_{L^2(\alpha)},
\]
\begin{equation}\label{eq-generalized-gradient-q-preconditioner}
\operatorname{PMO}_{\mathbb A,\alpha}(u)
=
-Q_\alpha^{-1}u,
\qquad
v_\alpha=-Q_\alpha^{-1}\Wgrad f(\alpha).
\end{equation}
\[
\mathbb A_p(\alpha,w)=\left(\int\norm{w}^p\d\alpha\right)^{2/p}.
\]
\begin{equation}\label{eq-wp-pmo}
\operatorname{PMO}_{\mathbb A_p,\alpha}(u)(x)
=
-\,U^{2-q}\norm{u(x)}^{q-2}u(x),
\end{equation}
The examples below are organized by the same dictionary. Concave-mobility flows use $\mathbb A_{\theta,\lambda}(\alpha,v)$ from Section~\ref{rem-generalized-bb}; spectral flows use $\mathbb A_\gamma(\alpha,v)$ from~\eqref{eq-spectral-tangent-action}; kernelized Benamou--Brenier flows use $\mathbb A_k(\alpha,v)=\norm{v}_{\RKHS_k^d}^2$ with a restricted RKHS tangent space.

\paragraph{General mobility flows.}
Let \(\lambda\) be Lebesgue measure, write \(\alpha=\rho\lambda\), and let \(u=\nabla(\delta f/\delta\rho)\). For the velocity action \(\mathbb A_{\theta,\lambda}\), the PMO introduced in Definition~\ref{def-generalized-action-steepest-direction} is the pointwise minimizer of

\(\int \dotp{u}{w}\,\rho\,\d x + \frac12\int \frac{\rho}{\theta(\rho)}\norm{w}^2\rho\,\d x,\)
\begin{equation}\label{eq-concave-mobility-pmo}
\operatorname{PMO}_{\mathbb A_{\theta,\lambda},\rho\lambda}(u)
=
-\frac{\theta(\rho)}{\rho}u.
\end{equation}
\begin{equation}\label{eq-general-mobility-gradient-flow}
\partial_t\rho_t
=
\nabla\!\cdot\!\left(\theta(\rho_t)\nabla\frac{\delta f}{\delta\rho}(\rho_t)\right).
\end{equation}
\(f(\rho)=\int U(\rho(x))\,\d x+\int V(x)\rho(x)\,\d x,\)
\[
\partial_t\rho
=
\nabla\!\cdot\!\left(\theta(\rho)\big(U''(\rho)\nabla\rho+\nabla V\big)\right).
\]
\paragraph{Dynamic spectral Wasserstein flows.}\label{sec-normalized-spectral-wasserstein-dynamics}
Specializing Definition~\ref{def-generalized-action-steepest-direction} to the spectral action \(\mathbb A_\gamma\), the descent direction associated with a field \(g\in L^2(\alpha;\RR^d)\) is simply \(\operatorname{PMO}_{\mathbb A_\gamma,\alpha}(g)\). The field \(g\) should be read as the Wasserstein gradient \(g=\Wgrad f(\alpha)=\nabla\delta f(\alpha)\).

\begin{prop}[Polar formula for the spectral direction]\label{prop-normalized-spectral-polar}
Let
\(S_\alpha(g)\eqdef\int g(x)g(x)^\top\d\alpha(x).\)
Assume that the inverse-trace problem admits a positive definite minimizer
\[
A_\alpha^\star
\in
\uargmin{A\in\mathcal B_\gamma\cap\mathbb S_{++}^d}
\tr\!\left(A^{-1}S_\alpha(g)\right).
\]
Then
\begin{equation}\label{eq-normalized-spectral-direction-explicit}
\operatorname{PMO}_{\mathbb A_\gamma,\alpha}(g)(x)
=
-(A_\alpha^\star)^{-1}g(x)
\end{equation}
is the spectral PMO direction. If the minimizer lies on the boundary of $\mathcal B_\gamma$, the same formula is understood by a limiting argument along positive definite minimizers, or equivalently by replacing the inverse by the corresponding inverse on the range of $S_\alpha(g)$.
\end{prop}
\begin{proof}
Using the polar formula $\gamma(M)=\sup_{A\in\mathcal B_\gamma}\tr(AM)$, the PMO objective~\eqref{eq-local-action-steepest-descent} with \(\mathbb A=\mathbb A_\gamma\) can be written as
\[
\sup_{A\in\mathcal B_\gamma}
\left\{
\int \dotp{g}{v}\d\alpha
+
\frac12\int v^\top A v\d\alpha
\right\}.
\]
For a fixed $A\succ0$, the quadratic lower bound
\(\int \dotp{g}{v}\d\alpha+\frac12\int v^\top A v\d\alpha\)
is minimized by $v_A=-A^{-1}g$, with value $-\frac12\tr(A^{-1}S_\alpha(g))$. Let $v_\star=-(A_\alpha^\star)^{-1}g$ and
\(M_\star=\int v_\star(x)v_\star(x)^\top\d\alpha(x) = (A_\alpha^\star)^{-1}S_\alpha(g)(A_\alpha^\star)^{-1}.\)
Since $A_\alpha^\star$ minimizes $A\mapsto\tr(A^{-1}S_\alpha(g))$ over the convex polar set, its first-order optimality condition gives
\(\tr(AM_\star)\leq \tr(A_\alpha^\star M_\star) \qquad \text{for all } A\in\mathcal B_\gamma.\)
Thus $A_\alpha^\star$ is active in the polar representation of $\gamma(M_\star)$, and
\(\gamma(M_\star)=\tr(A_\alpha^\star M_\star) = \tr((A_\alpha^\star)^{-1}S_\alpha(g)).\)
For any $v$, the polar formula gives \(\int \dotp{g}{v}\d\alpha+\frac12\mathbb A_\gamma(\alpha,v) \geq \int \dotp{g}{v}\d\alpha+\frac12\int v^\top A_\alpha^\star v\d\alpha,\)
and the right-hand side is minimized at $v_\star$. The activity identity above shows that equality holds at $v_\star$, hence $v_\star$ is the spectral PMO direction. The boundary case follows by approximation with positive definite matrices in the polar set.
\end{proof}
For the trace gauge, $\mathcal B_\gamma=\enscond{A}{0\preceq A\preceq\Id}$ and the minimizer is $A_\alpha^\star=\Id$. For the operator gauge $\gamma(M)=\lambda_{\max}(M)$, one has $\mathcal B_\gamma=\enscond{A\succeq0}{\tr(A)\leq1}$; if $S_\alpha(g)\succ0$, then

\[
A_\alpha^\star
=
\frac{S_\alpha(g)^{1/2}}{\tr(S_\alpha(g)^{1/2})},
\qquad
\operatorname{PMO}_{\mathbb A_\gamma,\alpha}(g)(x)
=
-\tr(S_\alpha(g)^{1/2})\,S_\alpha(g)^{-1/2}g(x).
\]
Proposition~\ref{prop-static-dynamic-spectral-wasserstein} identifies the static spectral distance $\Wass_\gamma$ with the length distance generated by $\mathbb A_\gamma$. The infinitesimal descent direction is therefore defined directly by the general PMO problem~\eqref{eq-local-action-steepest-descent} specialized to this action.

\begin{prop}[Formal normalized spectral gradient flow]\label{prop-normalized-spectral-gradient-flow}
Assume that $f$ has a smooth first variation and that the spectral PMO is attained along a smooth curve. The formal gradient flow generated by the dynamic spectral action $\mathbb A_\gamma$ solves
\begin{equation}\label{eq-normalized-spectral-flow}
\partial_t\alpha_t
+
\diverg\!\left(
\alpha_t\,\operatorname{PMO}_{\mathbb A_\gamma,\alpha_t}(\Wgrad f(\alpha_t))
\right)=0.
\end{equation}
Equivalently, if $A_t^\star$ solves the inverse-trace problem for
\(S_t=\int \Wgrad f(\alpha_t)(x)\Wgrad f(\alpha_t)(x)^\top\d\alpha_t(x),\)
then the velocity is \(v_t(x)=-(A_t^\star)^{-1}\Wgrad f(\alpha_t)(x).\)
Moreover, along smooth solutions, \(\frac{\d}{\d t}f(\alpha_t) = -\mathbb A_\gamma(\alpha_t,v_t).\)
If the minimizing movements~\eqref{eq-generalized-jko-step} for $d=\Wass_\gamma$ converge to a smooth curve with this minimal-action metric derivative, that limit satisfies the same equation.
\end{prop}
\begin{proof}
The action-based steepest-descent rule~\eqref{eq-generalized-gradient-from-action} gives
\(
v_t=\operatorname{PMO}_{\mathbb A_\gamma,\alpha_t}(\Wgrad f(\alpha_t))
\).
Proposition~\ref{prop-normalized-spectral-polar} gives the explicit inverse-trace formula.
Finally, the variation formula gives $\frac{\d}{\d t}f(\alpha_t)=\int\dotp{\Wgrad f(\alpha_t)}{v_t}\d\alpha_t$. Since $v_t$ minimizes a $2$-homogeneous penalized objective, the one-dimensional rescaling $s\mapsto sv_t$ has derivative zero at $s=1$, hence
\(\int\dotp{\Wgrad f(\alpha_t)}{v_t}\d\alpha_t + \mathbb A_\gamma(\alpha_t,v_t) =0,\)
which proves the dissipation identity.
\end{proof}
\subsection{Nonlocal Wasserstein Flows}
\label{sec-nonlocal-wasserstein-flows}
\label{sec-nonlocal-wasserstein-pdes}
\paragraph{Nonlocal Wasserstein flows and fractional PDEs.}
\begin{prop}[Entropy gradient flow for reversible jump kernels]\label{prop-nonlocal-entropy-gradient-flow}
Under the regularity and irreducibility assumptions used in Proposition~\ref{prop-nonlocal-distance-properties}, the Markov semigroup generated by the closure of
\(L\psi(x) = \int\bigl(\psi(y)-\psi(x)\bigr)K(x,\d y)\)
is the gradient flow, for the distance \(\mathcal W_K\) defined in \eqref{eq-nonlocal-wasserstein-distance}, of the relative entropy
\(\mathrm{Ent}_{\mathfrak m}(\alpha) = \int \rho\log\rho\,\d\mathfrak m, \qquad \alpha=\rho\mathfrak m.\)
Equivalently, in weak form, the entropy flow is
\begin{equation}\label{eq-nonlocal-entropy-flow}
\partial_t\rho_t=L\rho_t.
\end{equation}
When \(K\) is singular, the integral defining \(L\) is understood in the principal-value sense.
\end{prop}
\begin{proof}[Proof idea]
The key identity is the logarithmic-mean relation
\(\theta(a,b)(\log a-\log b)=a-b.\)
The first variation of \(\mathrm{Ent}_{\mathfrak m}\) is \(\log\rho+1\). With the sign convention of \eqref{eq-nonlocal-continuity-weak}, the steepest descent velocity is
\(v(x,y) = -\bar\nabla(\log\rho+1)(x,y) = \log\rho(x)-\log\rho(y).\)
Hence \(\theta(\rho(x),\rho(y))v(x,y)=\rho(x)-\rho(y)\). Substituting this in \eqref{eq-nonlocal-continuity-weak} and using the symmetry of \(\mathsf J\) gives
\(\frac{\d}{\d t}\int \varphi\rho\,\d\mathfrak m = \int \varphi(x)\int(\rho(y)-\rho(x))K(x,\d y)\,\mathfrak m(\d x),\)
which is the weak form of \(\partial_t\rho=L\rho\). The metric properties needed to justify this computation as a gradient-flow identity are the distance and geodesic results recalled in Proposition~\ref{prop-nonlocal-distance-properties}.
\end{proof}
\(K_s(x,\d y) = c_{d,s}\frac{\d y}{|x-y|^{d+s}}, \qquad 0<s<2.\)
\(L_s\psi(x) = \mathrm{p.v.}\int_{\RR^d} \bigl(\psi(y)-\psi(x)\bigr) \frac{c_{d,s}}{|x-y|^{d+s}}\,\d y = -(-\Delta)^{s/2}\psi(x),\)
\begin{equation}\label{eq-fractional-heat-nonlocal-wasserstein}
\partial_t\rho_t
=
-(-\Delta)^{s/2}\rho_t.
\end{equation}
More generally, when a target-dependent jump kernel \(K_\beta\) satisfies detailed balance with a desired invariant measure \(\beta=\rho_\beta\mathfrak m\), the same construction can be applied to \(r_t=\d\alpha_t/\d\beta\). The gradient flow of \(\KL(\alpha|\beta)\) is then \(\partial_t r_t=L_\beta r_t\).

\paragraph{Discrete Wasserstein flows on Markov chains.}
\(\operatorname{Ent}_\pi(\rho)\eqdef\sum_i\pi_i\rho_i\log\rho_i.\)
\(\frac{\rho^{k+1}-\rho^k}{\tau} =-\mathcal K_{\rho^{k+1}}\log\rho^{k+1}+o(1) =K\rho^{k+1}+o(1).\)
Here \((K\rho)_i=\sum_jK_{ij}(\rho_j-\rho_i)\). Under smoothness and strict positivity, $\rho^{k+1}=\rho^k+O(\tau)$, so the last expression also equals $K\rho^k+O(\tau)$.

\begin{prop}[Entropy gradient flow of a reversible Markov chain]
\label{prop-discrete-markov-entropy-gradient}
Let \(K\) be reversible with invariant law \(\pi\). The gradient flow of \(\operatorname{Ent}_\pi\) for the metric \(\mathcal W_K\) is the forward equation of the Markov chain,
\begin{equation}
\label{eq-discrete-markov-gradient-flow}
\dot\rho_i(t)=\sum_jK_{ij}\bigl(\rho_j(t)-\rho_i(t)\bigr).
\end{equation}
Equivalently, for the masses \(a_i(t)=\pi_i\rho_i(t)\), this is \(\dot a_i(t)=\sum_j(a_j(t)K_{ji}-a_i(t)K_{ij})\).
\end{prop}
\begin{proof}
The first variation of the entropy, with respect to the weighted pairing \(\sum_i\pi_i\xi_i\varphi_i\), is \(\log\rho_i+1\). Constants do not contribute to \(\mathcal K_\rho\), hence the metric gradient-flow equation is
\(\dot\rho=-\mathcal K_\rho\log\rho.\)
Using the identity
\(\theta(a,b)(\log a-\log b)=a-b,\)
one obtains, componentwise, \(\dot\rho_i =-\sum_jK_{ij}\theta(\rho_i,\rho_j)(\log\rho_i-\log\rho_j) = \sum_jK_{ij}(\rho_j-\rho_i),\)
which is the density form of the Kolmogorov forward equation. Multiplying by \(\pi_i\) and using detailed balance gives the equation for the probability masses.
\end{proof}
\subsection{Dynamic Unbalanced OT and WFR Flows}
\label{sec-dynamic-unbalanced-wfr-flows}
\label{sec-wfr-gradient-flows}
The generalized dynamic Wasserstein flows of Section~\ref{sec-generalized-dynamic-wasserstein-flows} are primarily mass-preserving: their tangent vectors are generated by continuity equations and are therefore velocity fields advecting the measure. Unbalanced transport changes the state space from probabilities to positive finite measures, where descent may both move and create or remove mass.

\begin{defn}[WFR minimizing movement]\label{def-wfr-jko-gradient-flow}
Let $f:\Mm_+(\RR^d)\to\RR\cup\{+\infty\}$ be an energy on positive finite measures and let $\WFR_\kappa$ be the Wasserstein--Fisher--Rao distance with growth scale $\kappa>0$, defined by the dynamic balance-equation formulation~\eqref{eq-dynamic-unbalanced-ot} and its static equivalence in Proposition~\ref{prop-static-dynamic-unbalanced}. The unbalanced JKO step is
\begin{equation}\label{eq-wfr-jko}
\alpha^{k+1}
\in
\uargmin{\alpha\in\Mm_+(\RR^d)}
\frac{1}{2\tau}\WFR_\kappa^2(\alpha^k,\alpha)+f(\alpha).
\end{equation}
The continuous-time limit, when it exists, is called the $\WFR_\kappa$ gradient flow of $f$.
\end{defn}
\(\partial_t\rho_t+\diverg(\rho_t v_t)=g_t\rho_t,\)
\begin{proposition}[Formal WFR gradient-flow PDE]\label{prop-wfr-gradient-flow-pde}
Assume that $f$ has a smooth first variation \(\phi_t=\delta f(\alpha_t)\) and that $\alpha_t=\rho_t\d x$ is smooth and positive. The formal $\WFR_\kappa$ gradient flow of $f$ is
\begin{equation}\label{eq-wfr-gradient-flow-pde}
\partial_t\rho_t
=
\diverg(\rho_t\nabla\phi_t)
-
\frac{1}{\kappa^2}\phi_t\rho_t.
\end{equation}
Along such a smooth solution,
\begin{equation}\label{eq-wfr-energy-dissipation}
\frac{\d}{\d t}f(\alpha_t)
=
-\int
\left(\norm{\nabla\phi_t}^2+\frac{1}{\kappa^2}\phi_t^2\right)\rho_t\,\d x
\leq 0.
\end{equation}
\end{proposition}
\begin{proof}
For an infinitesimal tangent pair $(v,g)$, the density variation is $\zeta=-\diverg(\rho v)+g\rho$. The first variation of $f$ is \(\int \phi\,\zeta = \int \dotp{\nabla\phi}{v}\rho\,\d x+\int \phi g\rho\,\d x,\)
after integration by parts. For the inner product
\(\dotp{(v,g)}{(\tilde v,\tilde g)}_\rho = \int \dotp{v}{\tilde v}\rho\,\d x+\kappa^2\int g\tilde g\rho\,\d x,\)
the Riesz representative is $(\nabla\phi,\phi/\kappa^2)$. Steepest descent therefore uses $(v,g)=(-\nabla\phi,-\phi/\kappa^2)$, which gives~\eqref{eq-wfr-gradient-flow-pde}.
Substituting this pair in the first-variation formula gives
\[
\frac{\d}{\d t}f(\alpha_t)
=
-\int
\left(\norm{\nabla\phi_t}^2+\frac{1}{\kappa^2}\phi_t^2\right)\rho_t\,\d x,
\]
which is~\eqref{eq-wfr-energy-dissipation}.
\end{proof}
\subsection{Conditional Wasserstein Training of Infinite ResNets}
\label{sec-conditional-wasserstein-resnets}
\paragraph{Finite-depth finite-width ResNets.}
\begin{equation}\label{eq-finite-resnet}
z_{r+1}
=
z_r+\tau G_{r,n_r}(z_r),
\qquad
G_{r,n_r}(z)
\eqdef
\frac1{n_r}\sum_{i=1}^{n_r}\psi(\theta_{r,i},z),
\qquad
z_0=z,\qquad r=0,\ldots,L-1.
\end{equation}
\(f_{L,n}((\theta_{r,i})_{r,i}) = \int \ell(z_L^{\theta}(z),y)\d\zeta(z,y).\)
\(\alpha_r = \frac1{n_r}\sum_{i=1}^{n_r}\delta_{\theta_{r,i}}, \qquad G_{r,n_r}(z) = \int_\Theta\psi(\theta,z)\d\alpha_r(\theta).\)
\(\alpha^{L,n}(\d s,\d\theta) = \frac1L\sum_{r=0}^{L-1}\delta_{s_r}(\d s)\alpha_r(\d\theta), \qquad s_r=r/L,\)
\paragraph{Finite-depth mean-field limit.}
\(z_{r+1} = z_r+\tau G_{\alpha_r}(z_r), \qquad z_0=z, \qquad G_{\alpha_r}(z) \eqdef \int_\Theta\psi(\theta,z)\d\alpha_r(\theta).\)
\(f_L(\alpha_0,\ldots,\alpha_{L-1}) = \int \ell(z_L^\alpha(z),y)\d\zeta(z,y).\)
\paragraph{Continuous-depth conditional geometry.}
\(\alpha(\d s,\d\theta)=\alpha_s(\d\theta)\d s,\)
or, equivalently, a probability measure on $[0,1]\times\Theta$ whose first marginal is Lebesgue measure. The general conditional Wasserstein distance is defined in Section~\ref{sec-conditional-wasserstein-distances}.

\begin{equation}\label{eq-resnet-depth-wasserstein}
\Wass_{2,\lambda}^2(\alpha,\beta)
=
\int_0^1\Wass_2^2(\alpha_s,\beta_s)\d s.
\end{equation}
\paragraph{Infinite-depth mean-field ResNet.}
\begin{equation}\label{eq-mean-field-resnet-forward}
\partial_s z_s
=
G_{\alpha_s}(z_s)
\eqdef
\int_\Theta \psi(\theta,z_s)\d\alpha_s(\theta),
\qquad z_0=z.
\end{equation}
\(f_{\mathrm{Res}}(\alpha) = \int \ell(z_1^\alpha(z),y)\d\zeta(z,y).\)
\paragraph{Conditional Wasserstein gradient flow.}
\[
\frac{\d}{\d\varepsilon}f(\alpha+\varepsilon\xi)\bigg|_{\varepsilon=0}
=
\int_S\int_\Theta \frac{\delta f}{\delta\alpha_s}(\alpha)(\theta)\d\xi_s(\theta)\d\lambda(s).
\]
\begin{equation}\label{eq-conditional-resnet-gradient-flow}
\partial_t\alpha_{t,s}
+\diverg_\theta(\alpha_{t,s}v_{t,s})=0,
\qquad
v_{t,s}(\theta)
=
-\nabla_\theta\frac{\delta f}{\delta\alpha_s}(\alpha_t)(\theta),
\end{equation}
for $\lambda$-a.e. condition $s$, whenever the first variation is differentiable in the parameter variable. In the ResNet case one takes $f=f_{\mathrm{Res}}$ and $S=[0,1]$.

\paragraph{Particle discretization and layerwise matching.}
\(\alpha_r=\frac1{n_r}\sum_{i=1}^{n_r}\delta_{\theta_{r,i}}, \qquad \beta_r=\frac1{n_r}\sum_{i=1}^{n_r}\delta_{\theta'_{r,i}},\)
\(\Wass_{2,\lambda_L}^2(\alpha^{L,n},\beta^{L,n}) = \frac1L\sum_{r=0}^{L-1} \Wass_2^2(\alpha_r,\beta_r).\)
\(\frac1L\sum_{r=0}^{L-1}\frac1{n_r} \sum_{i=1}^{n_r}\norm{\theta_{r,i}-\theta'_{r,i}}^2,\)
\subsection{Second-Order Momentum Flows}
\label{sec-second-order-momentum-flows}
\paragraph{Finite-dimensional momentum.}
Let $F:\RR^{N}\to\RR$ be a smooth finite-dimensional objective. With step size $h>0$ and damping $\lambda_{\rm fric}>0$, one convenient velocity form is

\begin{equation}\label{eq-polyak-momentum}
S^{k+1}=(1-\lambda_{\rm fric} h)S^k-h\nabla F(X^k),
\qquad
X^{k+1}=X^k+hS^{k+1}.
\end{equation}
If $S^k$ approximates $\dot X(kh)$, this is a semi-implicit, or symplectic-Euler, discretization of the first-order system: the velocity update is explicit, while the position uses the newly updated velocity. The limiting system is

\begin{equation}\label{eq-heavy-ball-first-order-system}
\dot X(t)=S(t),
\qquad
\dot S(t)=-\lambda_{\rm fric} S(t)-\nabla F(X(t)),
\end{equation}
\begin{equation}\label{eq-heavy-ball-ode}
\ddot X(t)+\lambda_{\rm fric}\dot X(t)+\nabla F(X(t))=0.
\end{equation}
\begin{equation}\label{eq-nesterov-momentum}
Y^k=X^k+\theta_k(X^k-X^{k-1}),
\qquad
X^{k+1}=Y^k-h\nabla F(Y^k),
\end{equation}
\begin{equation}\label{eq-nesterov-ode}
\ddot X(t)+\frac{r}{t}\dot X(t)+\nabla F(X(t))=0.
\end{equation}
\paragraph{Empirical Wasserstein lift.}
\begin{equation}\label{eq-empirical-momentum-lift}
\al_X \eqdef \frac1n\sum_{i=1}^n \delta_{x_i},
\qquad
F(X) \eqdef f(\al_X).
\end{equation}
\(\norm{\dot X}_n^2 \eqdef \frac1n\sum_{i=1}^n\norm{\dot x_i}^2,\)
\[
\left.\frac{\d}{\d\varepsilon}\right|_{\varepsilon=0}
F(x_1,\ldots,x_i+\varepsilon\xi,\ldots,x_n)
=
\frac1n\dotp{\nabla\delta f(\al_X)(x_i)}{\xi}.
\]
Thus, if $\nabla^{(n)}F$ denotes the gradient for the metric $\norm{\cdot}_n$, then

\begin{equation}\label{eq-empirical-lift-gradient}
(\nabla^{(n)}F(X))_i
=
n\nabla_{x_i}F(X)
=
\nabla\delta f(\al_X)(x_i)
=
\Wgrad f(\al_X)(x_i).
\end{equation}
\begin{equation}\label{eq-momentum-wasserstein-velocity}
v_\al(x)\eqdef -\Wgrad f(\al)(x)=-\nabla\delta f(\al)(x).
\end{equation}
\begin{equation}\label{eq-wasserstein-momentum-particles}
\dot x_i(t)=s_i(t),
\qquad
\dot s_i(t)=-\lambda_{\rm fric} s_i(t)+v_{\al_t}(x_i(t)),
\qquad
\al_t=\frac1n\sum_{i=1}^n\delta_{x_i(t)}.
\end{equation}
\(\ddot x_i(t)+\lambda_{\rm fric}\dot x_i(t) = v_{\al_t}(x_i(t)) = -\Wgrad f(\al_t)(x_i(t)).\)
\begin{prop}[Discrete mechanical energy dissipation]\label{prop-second-order-energy-dissipation}
Assume that $f$ and the particle solution are smooth enough to justify differentiating the identities above along \eqref{eq-wasserstein-momentum-particles}. Define
\(\mathcal H_n(X,S) \eqdef F(X)+\frac1{2n}\sum_{i=1}^n\norm{s_i}^2, \qquad S=(s_1,\ldots,s_n).\)
Then along \eqref{eq-wasserstein-momentum-particles},
\begin{equation}\label{eq-second-order-energy-dissipation}
\frac{\d}{\d t}\mathcal H_n(X(t),S(t))
=
-\frac{\lambda_{\rm fric}}{n}\sum_{i=1}^n\norm{s_i(t)}^2
\leq 0.
\end{equation}
\end{prop}
\begin{proof}
By \eqref{eq-empirical-lift-gradient} and the definition $v_\al=-\Wgrad f(\al)$, one has $\nabla_{x_i}F(X)=-v_{\al_X}(x_i)/n$. Hence
\(\frac{\d}{\d t}F(X(t)) = \sum_{i=1}^n \dotp{\nabla_{x_i}F(X(t))}{s_i(t)} = -\frac1n\sum_{i=1}^n\dotp{v_{\al_t}(x_i(t))}{s_i(t)}.\)
On the other hand, \(\frac{\d}{\d t}\frac1{2n}\sum_i\norm{s_i(t)}^2 = \frac1n\sum_i\dotp{s_i(t)}{-\lambda_{\rm fric} s_i(t)+v_{\al_t}(x_i(t))}.\)
The force terms cancel, leaving \eqref{eq-second-order-energy-dissipation}.
\end{proof}
\paragraph{Phase-space formulation.}
\begin{prop}[Phase-space Liouville formulation]\label{prop-second-order-liouville}
Let $v_\al$ be smooth enough that \eqref{eq-wasserstein-momentum-particles} has a classical solution. Define the empirical phase-space measure
\(\eta_t^n \eqdef \frac1n\sum_{i=1}^n\delta_{(x_i(t),s_i(t))} \in \Pp(\RR^d_x\times\RR^d_s), \qquad \al_t^n=(\pi_x)_\sharp\eta_t^n.\)
Then $\eta_t^n$ solves, in the sense of distributions,
\begin{equation}\label{eq-second-order-liouville-empirical}
\partial_t\eta_t^n
+
\nabla_x\cdot(s\,\eta_t^n)
+
\nabla_s\cdot\bigl((-\lambda_{\rm fric} s+v_{\al_t^n}(x))\,\eta_t^n\bigr)
=0.
\end{equation}
If, as $n\to\infty$, the empirical phase-space laws converge to a smooth density $\eta_t(x,s)$ and the nonlinear fields converge accordingly, the limiting kinetic equation is
\begin{equation}\label{eq-second-order-liouville-density}
\partial_t\eta_t
+
\nabla_x\cdot(s\,\eta_t)
+
\nabla_s\cdot\bigl((-\lambda_{\rm fric} s+v_{\al_t}(x))\,\eta_t\bigr)
=0,
\qquad
\al_t=(\pi_x)_\sharp\eta_t.
\end{equation}
\end{prop}
\begin{proof}
Let $\varphi\in C_c^\infty(\RR^d_x\times\RR^d_s)$. Along the particle system,
\[
\frac{\d}{\d t}
\int \varphi\,\d\eta_t^n
=
\frac1n\sum_{i=1}^n
\left[
\dotp{\nabla_x\varphi(x_i,s_i)}{s_i}
+
\dotp{\nabla_s\varphi(x_i,s_i)}{-\lambda_{\rm fric} s_i+v_{\al_t^n}(x_i)}
\right].
\]
Rewriting the right-hand side as an integral against $\eta_t^n$ and integrating by parts gives \eqref{eq-second-order-liouville-empirical}. Passing formally to the mean-field limit gives \eqref{eq-second-order-liouville-density}.
\end{proof}
\(k(x,y)=-\norm{x-y}, \qquad f(\al)=\MMD_k^2(\al,\beta) = \iint k(x,y)\,\d(\al-\beta)(x)\d(\al-\beta)(y).\)
\(v_\al(x) = 2\int \frac{x-y}{\norm{x-y}}\,\d(\al-\beta)(y).\)
\paragraph{Entropy without an empirical energy.}
For the numerical illustration, we add the quadratic confinement $V(x)=\norm{x}^2/2$. In one dimension the overdamped equation becomes the Ornstein--Uhlenbeck Fokker--Planck equation

\(\partial_t\rho_t=\partial_{xx}\rho_t+\partial_x(x\rho_t),\)
whose stationary law is the centered Gaussian $\Gaussian(0,1)$. The Newton lift is discretized directly in phase space: for a density $\eta_t(x,s)$ with spatial marginal $\rho_t(x)=\int\eta_t(x,s)\,\d s$, it reads

\(\partial_t\eta_t+\partial_x(s\eta_t)+\partial_s\bigl((-\partial_x\log\rho_t(x)-x)\eta_t\bigr)=0.\)